% ek_rapor.tex — "Revizyon kontrolleri" raporuna gelen 47 soruya cevaplar (ek rapor)
% Derleme: xelatex ek_rapor.tex (iki kez)
\documentclass[11pt,a4paper]{article}
\usepackage{fontspec}
\usepackage{polyglossia}
\setdefaultlanguage{turkish}
\setotherlanguage{english}
\usepackage{amsmath,amssymb}
\usepackage{graphicx}
\usepackage{booktabs}
\usepackage{array}
\usepackage{tabularx}
\usepackage[margin=2.3cm]{geometry}
\usepackage{xcolor}
\usepackage{xurl}
\usepackage{hyperref}
\usepackage{caption}
\usepackage{float}
\usepackage{enumitem}
\hypersetup{colorlinks=true,linkcolor=blue!50!black,urlcolor=blue!50!black,citecolor=blue!50!black}
\captionsetup{font=small,labelfont=bf}
\setlength{\parskip}{0.45em}
\graphicspath{{sonuclar/}}
\newcommand{\sekil}[3][0.95]{\begin{figure}[H]\centering\includegraphics[width=#1\textwidth]{#2}\caption{#3}\end{figure}}
\setlength{\parindent}{0pt}

\newcommand{\kod}[1]{{\XeTeXinterchartokenstate=0\texttt{#1}}}
\newcommand{\BEL}{\textsc{[belgelenmiş]}}
\newcommand{\YOR}{\textsc{[yorum]}}
\newcommand{\HIP}{\textsc{[hipotez]}}
\newcommand{\LIT}{\textsc{[literatür]}}

% Her soru: numara, alıntı kutusu, cevap
\newcounter{soru}
\newcommand{\soru}[2]{%
  \par\medskip\noindent\textbf{Soru #1.}\enspace\textit{#2}\par\nopagebreak\smallskip}
\newcommand{\alinti}[1]{\begin{quote}\small\color{black!70}``#1''\end{quote}}

\title{Ek Rapor: Revizyon Kontrolleri Üzerine Sorulara Cevaplar}
\author{Hazırlayan: Claude (Anthropic yapay zekâ asistanı)\\
\small ``Ters Yönde Dönen Kapalı Yörüngelerle Simetrik Kaçıklık İzleme:\\ \small PRAB Revizyonu Öncesi
Geçerlilik Kontrolleri'' raporuna ektir}
\date{8 Kasım 2026}

% otomatik: ek_rapor_sonuc_uret.py
\newcommand{\Ceqnine}{0.014 [0.014, 0.014] (n=1)}
\newcommand{\Ceqnf}{0.0078 [0.0078, 0.0078] (n=1)}
\newcommand{\Cetnine}{0.00084 [0.00084, 0.00084] (n=1)}
\newcommand{\Cetnf}{0.01 [0.01, 0.01] (n=1)}
\newcommand{\RhoBirEbir}{13 [13, 13] & 1/1}
\newcommand{\RhoBirEiki}{0.11 [0.11, 0.11] & 0/1}
\newcommand{\RhoBirEuc}{bekleniyor & --}
\newcommand{\RhoBirEdort}{bekleniyor & --}
\newcommand{\EdSifirKuv}{--}
\newcommand{\EdSifirC}{bekleniyor & -- & --}
\newcommand{\EdMerkezKuv}{--}
\newcommand{\EdMerkezC}{bekleniyor & -- & --}
\newcommand{\RhoYorum}{\textit{(C++ sonuçları bekleniyor.)}}
\newcommand{\RhoBesYorum}{\textit{(C++ sonuçları bekleniyor.)}}
\newcommand{\EdYorum}{\textit{(C++ sonuçları bekleniyor.)}}

\begin{document}
\sloppy
\maketitle

\begin{abstract}
\noindent Bu ek rapor, ana rapor üzerine iletilen 47 soruyu, sorulardaki numaralarla ve sırayla cevaplar.
Cevapların bir kısmı için yeni hesaplar yapıldı: (i)~fırlatma noktasına bağımlılığın gerçek nedeni
(sert kenarlı saptırıcılarda $|p|$'nin hücreden hücreye değişmesi; ana rapordaki ``potansiyel enerji''
açıklaması yanlıştı ve düzeltildi), (ii)~$\rho$ yanlış bilinirse ne olur, (iii)~$\rho=0.5$'te sahte EDM,
(iv)~elektrostatik düzelticinin (Ed) yeni prosedürle sahte EDM'si, (v)~TSVD eşiğinin doğru kaçıklığı
bilmeden seçilmesi. Literatür değerleri (BPM kararlılığı, BPM gürültüsü, yüksek gerilim kararlılığı)
kaynaklarıyla verildi. En önemli yeni bulgu olumsuzdur: saptırıcı başına alan kararlılığı için
varsaydığımız $\sim3\times10^{-7}$ düzeyi, bugünkü yüksek gerilim kaynaklarının
($10^{-5}$--$10^{-4}$) çok altındadır; bu, yöntemin en zor tasarım gereksinimidir.
\end{abstract}

\tableofcontents
\newpage

%======================================================================
\section{Bu belge nasıl okunmalı}
%======================================================================
Her soru, soru listesindeki numarasıyla başlar; ardından ana rapordaki alıntı ve cevap gelir. Ana raporun
etiketleri burada da kullanılıyor: \BEL\ bu çalışmada çalıştırılmış bir hesabın çıktısı, \YOR\ belgelenmiş
sonuçlardan benim çıkarımım, \HIP\ test edilmemiş tahmin. Ek olarak \LIT\ bir yayından alınmış değer
demektir; kaynak listesi sondadır (Kaynaklar bölümü).

Yeni hesapların hepsi \kod{claude/brave-allen-92wn61} dalında, \kod{revizyon/} altındadır:
\kod{tests/teshis\_p\_azimut.py} (soru 9), \kod{ek\_testler.py} (soru 15, 16, 22),
\kod{ek\_rcond\_secimi.py} (soru 14). Sahte EDM'ler donmuş spin prosedürüyle, kısaltılmış ayarla
(sayısal taban $\approx0.008\times$hedef) ölçüldü. $C=\tfrac12|f_{\rm CW}-f_{\rm CCW}|$, hedefin katı.

\paragraph{Ana raporda yapılan düzeltmeler.} Bu sorular sırasında ana raporda üç yer düzeltildi:
(1)~Bölüm 3.2'deki ``Neden?'' listesi (soru 9), (2)~$E_z\to E_y$ gösterimi (soru 18), (3)~yeni ölçüm
aracının 1. adımına ``enerji korunumlu fırlatma yalnız başlangıç tahminidir'' notu (soru 8). Ayrıca
Adım~6 gürültü modelinde tur frekansı 270~kHz alınmıştı; doğrusu 224~kHz ($T_{\rm dev}=4.4577~\mu$s).
Etkisi ihmal edilebilir: $\sigma_n$ 10.31 yerine 10.37~nm (soru 35).

%======================================================================
\section{Temel büyüklükler (ana rapor 1.2)}
%======================================================================

\soru{1}{CW ve CCW matrisleri neden ayrı? Aralarındaki fark ne? Simetrik ve antisimetrik bileşen bu iki
matriste nasıl farklı görünüyor?}

Tepki matrisi $R$, kuadrupol kaçıklığından BPM okumasına giden doğrusal haritadır. Kaçık bir kuadrupol,
merkezinden geçmeyen demete bir dipol alanı uygular: $\Delta B = g\,d$ ($g$ gradyan, $d$ kaçıklık). Bu
alanın demete verdiği tekme $\theta=q\,v\times\Delta B$'den gelir ve hız yönüne bağlıdır. Demet yönü
çevrilince tekme işaret değiştirir. Sonra bu tekme halka boyunca bir kapalı yörünge bozulmasına dönüşür;
bu ikinci adım (betatron hareketi) iki yönde aynı yapıdadır.

Sonuç: $R_{\rm CW}\approx -R_{\rm CCW}$ olması beklenir. Ölçülen sütun korelasyonu iki düzlemde de
yaklaşık $-0.95$'tir~\BEL\ (test paketi, \kod{test\_cw\_ccw\_matris\_isaret\_yapisi}). Tam $-1$ olmamasının
nedeni, saptırıcıların elektrik alanının yönle işaret değiştirmemesidir. Demeti büken kuvvetin bir kısmı
elektrik, bir kısmı manyetik olduğundan iki yönün optiği tam ayna görüntüsü değildir.

Simetrik ve antisimetrik bileşenin iki matristeki görünüşü şöyle. Sayılar Frobenius normlarıdır; $P$
simetrik ya da antisimetrik izdüşüm, ``toplam kanalı'' $\tfrac12\|(R_{\rm CCW}+R_{\rm CW})P\|$, ``fark kanalı''
$\tfrac12\|(R_{\rm CCW}-R_{\rm CW})P\|$~\BEL:
\begin{center}\small
\begin{tabular}{lcccc}
\toprule
& $\|R_{\rm CCW}P\|$ & $\|R_{\rm CW}P\|$ & toplam kanalı & fark kanalı \\
\midrule
x, simetrik & 8.0 & 8.1 & 5.2 & 6.2 \\
x, antisimetrik & 33.6 & 33.6 & 5.1 & 33.2 \\
y, simetrik & 9.5 & 9.5 & 6.9 & 6.5 \\
y, antisimetrik & 45.8 & 45.8 & 6.8 & 45.3 \\
\bottomrule
\end{tabular}
\end{center}
\sekil{ek_sekil_s1_kanallar.pdf}{Simetrik ve antisimetrik kaçıklığın toplam (CW+CCW) ve fark (CW$-$CCW) kanallarındaki yörünge tepkisi. Antisimetrik sinyal neredeyse tamamen fark kanalında; toplam kanalında ikisi eşit büyüklükte.}

Okunuşu: antisimetrik kaçıklık büyük bir yörünge üretir (norm $\sim$34--46) ve bunun neredeyse tamamı
fark kanalındadır. Simetrik kaçıklık 4--5 kat küçük bir yörünge üretir ve bu yörüngenin yarısına yakını
toplam kanalındadır. Toplam kanalında simetrik ve antisimetrik eşit büyüklükte görünür (5.2 ve 5.1).
Antisimetriği ayırt eden fark kanalıdır. Makalenin fikri buradan gelir: tek demette simetrik sinyal,
çok daha büyük antisimetrik sinyalin altında kalır; iki demetle fark kanalı antisimetriği ayrı ölçer ve
geri kalanı simetriğe bırakır~\YOR.

%======================================================================
\section{Adım 0: Kodda ne varsayılmış?}
%======================================================================

\soru{2}{``Düzeltme, kestirilen kaçıklığın doğrudan kuadrupol merkezinden çıkarılmasıyla yapılıyor.''
Bu, simülasyonda tam olarak ne demek? Neden kuadrupolün içindeki manyetik düzelticiye eşdeğer? Gerçek bir
makinede karşılığı ne?}

Simülasyonda her kuadrupolün merkezi \kod{quad\_dx}, \kod{quad\_dy} dizileriyle verilir. Gerçek
kaçıklık $d$ ise kestirim $\hat d$ olsun. ``Düzeltme'' bu dizilere $d-\hat d$ yazmaktır: C++ izleyici,
kuadrupolü sanki merkezi $d-\hat d$ kadar kaymış gibi izler.

Neden manyetik düzelticiye eşdeğer: kuadrupolün alanı $B=g\,(y-d)$ biçimindedir (düşey düzlem için). Aynı
kuadrupolün içine düzgün bir dipol alanı $\Delta B$ eklenirse alan $g\,(y-d)+\Delta B=g\,(y-[d-\Delta B/g])$
olur. Yani düzgün dipol alanı, kuadrupolün manyetik merkezini $\Delta B/g$ kadar kaydırmakla
\emph{birebir aynıdır}; hem yörünge hem spin için. $\Delta B=g\,\hat d$ seçilirse etkin merkez $d-\hat d$
olur~\YOR.

Gerçek makinede iki karşılığı vardır~\YOR: (a)~kuadrupolün kutuplarına sarılmış dipol trim sargıları
(birçok halkada ``kombine işlevli'' düzeltici olarak kullanılır); (b)~kuadrupolü mekanik olarak taşıyan
hassas konumlandırıcılar (ör.\ kam tahrikli ``magnet mover''lar). (a) sürekli ve hızlı, (b) yavaş ama
alan kalitesini bozmaz. Bu ayrıntı makalede yazılmalıdır (soru 3).

\soru{3}{Düzeltici türünün yazılmaması neden önemli bir eksik? Sonuç düzelticinin türüne ne kadar bağlı?}

Çok bağlı. Ana rapor Adım~3'te aynı kestirimle dört düzeltici türü denendi. Kuadrupol içi düzeltici (Mq)
$C\approx10^{-3}$ verirken saptırıcı konumundaki manyetik düzeltici 0.04--55, elektrostatik düzeltici
1.5--14 verdi~\BEL\ (eski prosedür; Ed için yeni ölçüm soru 22'de). Yani doğru düzeltici türü, sonucu
$10^3$--$10^4$ kat değiştiriyor. Bir okuyucu ``yörüngeyi düzelttik'' cümlesini standart yörünge
düzelticileriyle (saptırıcı yanındaki dipoller) yorumlarsa, makalenin sayılarını yeniden üretemez. Bu
yüzden düzeltici türü, sonucun kendisi kadar önemlidir~\YOR.

\soru{4}{$S=[R_{\rm CCW};R_{\rm CW}]$, \kod{rcond}~$=0.01$ ile TSVD. Bu fit adım adım nasıl kuruluyor?
İki matrisi alt alta koymak ne kazandırıyor? \kod{rcond}~$=0.01$ ne demek ve neden seçilmiş?}

\textbf{Kuruluş.}
\begin{enumerate}[nosep]
\item İki demetin 48'er BPM okuması alt alta yazılır: $y=[y_{\rm CCW};y_{\rm CW}]$ (96 sayı).
\item Model: $y=S\,d+n$, $S=[R_{\rm CCW};R_{\rm CW}]$ ($96\times48$), $n$ gürültü.
\item $S$'nin tekil değer ayrışımı alınır: $S=U\Sigma V^{T}$. Tekil değerler $s_1\ge s_2\ge\dots\ge s_{48}$.
\item $s_i<\kod{rcond}\cdot s_1$ olan doğrultular atılır; kalanlarla $\hat d=\sum_{i\,\rm kalan} v_i\,(u_i^{T}y)/s_i$.
\end{enumerate}
Atılan doğrultular, verinin çok zayıf gördüğü kaçıklık desenleridir. Onları tutmak, gürültüyü $1/s_i$
kadar büyütür.

\textbf{Alt alta koymanın kazancı.} Tek demette (yalnız $R_{\rm CCW}$) koşul sayısı x'te 135, y'de 228;
\kod{rcond}~$=0.01$ 48 doğrultunun 17'sini (x) ve 24'ünü (y) atar. İki demette koşul sayısı 73 ve 118'e
iner; aynı eşik x'te hiç, y'de 8 doğrultu atar~\BEL. Atılan doğrultular ağırlıklı olarak simetrik
desenlerdir (soru 1: simetrik yörünge küçük). Yani tek demet, \kod{rcond}~$=0.01$ ile simetrik
bileşenin büyük kısmını ``bilinmiyor'' sayıp sıfır kabul eder; iki demet bunları tutar~\YOR.

\sekil{ek_sekil_s2_tekil.pdf}{Tek demet ($R_{\rm CCW}$) ve iki demet ($[R_{\rm CCW};R_{\rm CW}]$) matrislerinin normalize tekil değerleri. Kesikli çizginin altındaki doğrultular \kod{rcond}~$=0.01$ ile atılır.}
\textbf{0.01 neden?} Makalede gerekçe yok; büyük olasılıkla sezgisel bir seçimdir. Ana rapor Adım~1
gösterdi ki bu seçim tek demeti haksız yere kötü gösteriyor: tek demet uygun eşikle (0.003) simetrik
artığı 3.87~$\mu$m yerine 0.051~$\mu$m'ye indiriyor~\BEL. Eşiğin gerçek kaçıklığı bilmeden nasıl seçileceği
soru 14'te.

\soru{5}{100~$\mu$m BPM ofsetiyle sonuç neden bu kadar kötüleşiyor? Bu, yöntem için ne anlama geliyor?
Ofset nasıl ele alınmalı?}

BPM ofseti, BPM'in elektriksel merkezinin bilinmeyen sabit kaymasıdır. Okumaya $b$ olarak eklenir ve
\emph{iki demette aynıdır}, çünkü aynı BPM'dir. Yani ofset tamamen toplam kanalına düşer. Soru 1'e göre
simetrik kaçıklığın sinyali de büyük ölçüde toplam kanalındadır. İki demetli fit, 100~$\mu$m'lik ofseti
ayırt edemez ve onu simetrik kaçıklık sanar. Simetrik yörünge tepkisi küçük olduğundan (norm $\sim$5--7),
100~$\mu$m ofset onlarca mikrometrelik sahte bir simetrik kaçıklık kestirimine dönüşür. Düzeltme bunu
kuadrupollere uygulayınca sahte EDM hedefin $\sim$200 katına çıkar~\YOR.

Anlamı: yöntem \emph{mutlak} kaçıklık ölçümü olarak kullanılamaz. Ancak \emph{sürüklenme izleyicisi}
olarak kullanılabilir: BBA (demet tabanlı hizalama; demetin kuadrupol merkezinden geçtiği BPM okumasının
gradyan modülasyonuyla bulunması) sonrası okuma bir referans olarak kaydedilir ve sonraki ölçümler bu
referanstan \emph{farkla} değerlendirilir. Sabit ofset farkta sadeleşir. Geriye yalnızca ofsetin
\emph{değişimi} kalır; ana rapordaki $\delta$ budur (T2, Adım 6). Makale bu ayrımı açıkça yazmalı:
yöntem ``BBA'dan sonra, BBA referansına göre sürüklenme izleme'' yöntemidir~\YOR.

%======================================================================
\section{Fırlatma enerjisi sorunu ve yeni ölçüm aracı (ana rapor 3.2--3.3)}
%======================================================================

\soru{6}{Sahte EDM neden enerji sapmasına bu kadar duyarlı? Fiziksel mekanizma ne? Bu iki sayı nasıl
ölçüldü?}

\textbf{Nasıl ölçüldü.} $\sigma=10~\mu$m, seed~0 desende, kapalı yörünge üzerinden fırlatılan parçacığın
momentumu $(1+\delta)$ ile çarpılarak $\delta=-10^{-6},-10^{-7},0,+10^{-7},+10^{-6}$ için beş ayrı izleme
yapıldı. Her birinde $f=\mathrm{d}S_y/\mathrm{d}t$ (model fit) ve yatay presesyon hızı
$\omega_h=\mathrm{d}/\mathrm{d}t\,\mathrm{atan2}(S_x,S_z)$ ölçüldü. İkisi de $\delta$'ya doğrusal çıktı;
eğimler $\mathrm{d}f/\mathrm{d}\delta\approx-2.9$~rad/s (yani $\delta=10^{-7}$ başına $\approx290\times$hedef)
ve $\mathrm{d}\omega_h/\mathrm{d}\delta\approx-5.24\times10^{6}$~rad/s~\BEL\ (\kod{tests/teshis\_enerji\_ABC.py}).

\sekil{ek_sekil_s3_delta.pdf}{Beş noktalı enerji taraması ($\sigma=10~\mu$m, seed~0). İkisi de doğrusal; elmas, $\omega_h=0$ olan donmuş spin noktası.}
\textbf{İki eğim çok farklı nitelikte~\BEL.} Donmuş spin aracının $\sigma=2.5$, 5, 10~$\mu$m koşumlarından:
\begin{itemize}[nosep]
\item $\mathrm{d}\omega_h/\mathrm{d}\delta=-5.238\times10^{6}$~rad/s, kaçıklıktan \emph{bağımsız}
(9 koşumun hepsinde 4 basamak aynı). Bu, halkanın kendi özelliğidir: sihirli momentumdan sapan bir
protonun spini elektrik alanında momentumdan hızlı ya da yavaş döner.
\item $\mathrm{d}f/\mathrm{d}\delta$ kaçıklıkla \emph{doğrusal} ölçeklenir: seed~0 için $\sigma=2.5$, 5,
10~$\mu$m'de $-0.72$, $-1.45$, $-2.91$~rad/s.
\end{itemize}

\textbf{Mekanizma~\YOR.} Donmuş spin noktasında spin, yatay düzlemde momentumu izler. Kapalı yörüngenin
düşey eğimleri spini radyal eksen etrafında küçük açılarla döndürür; bir tur boyunca bu dönmelerin
birinci mertebe toplamı sıfıra kapanır (kapalı yörünge koşulu). Geriye ikinci mertebe, $\sigma^2$ ile
ölçeklenen geometrik faz kalır; bu, aradığımız sahte EDM'dir. $\delta\ne0$ iken spin her turda yatayda
$\omega_h$ hızıyla döner. Bu ek dönme, düşey dönmelerin birinci mertebedeki sıfıra kapanmasını bozar ve
$\omega_h\times(\text{kaçıklıkla doğrusal düşey dönme})$ mertebesinde bir terim çıkar. Bu terim
$\sigma^1\delta$ ile ölçeklenir; ölçülen $\mathrm{d}f/\mathrm{d}\delta\propto\sigma$ ile tutarlıdır.
Duyarlılığın büyüklüğü buradan gelir: $\delta=10^{-7}$ için $\omega_h\approx0.5$~rad/s, hedef ise
$10^{-9}$~rad/s. Yatay presesyonun çok küçük bir kısmının düşeye ``sızması'' bile hedefi aşar.

Kaba analitik kontrol: sihirli olmayan momentumda tek bir saptırıcı içindeki göreli presesyon formülü test
paketinde doğrulandı. Bu formülün yalnız $G$ terimi ve saptırıcı doluluk oranıyla verdiği
$\mathrm{d}\omega_h/\mathrm{d}\delta\approx2\times10^{6}$~rad/s, ölçülenle aynı mertebededir; kalan
$\sim2.5$ kat farkı yörüngenin enerjiyle değişmesine (dispersiyon) bağlıyorum~\HIP.

\soru{7}{``Deneyde sabit tutulan şey $\omega_h=0$ koşuludur.'' Bunu açar mısın? Gerçek deneyde bu geri
besleme nasıl çalışıyor ve neden simülasyonda doğru referans noktası?}

Polarimetre, spinin yatay düzlemdeki yönünü (fazını) sürekli ölçer. Spin donmuşsa bu faz sabittir; değilse
$\omega_h$ hızıyla döner. JEDI işbirliği COSY'de spin presesyon fazını ölçüp RF frekansını geri beslemeyle
ayarlayarak fazı kilitlemeyi göstermiştir~\LIT~\cite{jedi_phaselock}. Donmuş spin halkasında aynı ilke, demet
enerjisini (ya da saptırıcı alanını) $\omega_h=0$ tutacak biçimde ayarlamaktır. Yani deneyde ``doğru
enerji'' tanımı, $\omega_h=0$'dır.

Neden simülasyonda doğru referans: Sahte EDM $\delta$'ya çok duyarlı olduğundan (soru 6), simülasyonda hangi
enerjiyi seçtiğimiz sonucu belirler. Deneyle karşılaştırılabilir tek seçim, deneyin kendisinin sabit
tuttuğu koşuldur. Momentumu ya da toplam enerjiyi sabitlemek, deneyde gerçekleşmeyen bir çalışma noktası
seçmek olur~\YOR. Ek bir avantaj: donmuş spin yörüngesi tektir; hangi noktadan bulunursa bulunsun aynı
yörüngedir (soru 9).

Geri beslemenin sonlu hassasiyeti: kalan hata $\Delta\omega_h$, $f$'ye
$(\mathrm{d}f/\mathrm{d}\delta)/(\mathrm{d}\omega_h/\mathrm{d}\delta)\cdot\Delta\omega_h$ kadar geçer. Bu oran
$\sigma=10~\mu$m'de $5.5\times10^{-7}$; düzeltme sonrası artık kaçıklıkta (0.01--0.3~$\mu$m) 30--1000 kat
küçüktür~\BEL\ (oran $\propto\sigma$, soru 6). Örneğin $\Delta\omega_h=10^{-3}$~rad/s ve 0.3~$\mu$m artık
için katkı $\sim2\times10^{-11}$~rad/s, yani hedefin \%2'si~\HIP\ ($\Delta\omega_h$ varsayımı).

\soru{8}{Yeni ölçüm aracının prosedürünü adım adım açıklar mısın? ``Enerji korunumlu fırlatma
($\delta=0$)'' tam olarak ne demek? İki noktalı doğrusal ara değerleme neden yeterli?}

\textbf{Adımlar.}
\begin{enumerate}[nosep]
\item \emph{Kapalı yörüngeyi bul.} 4 boyutlu ($x,x',y,y'$) arama: betatron salınımının genliğini en küçük
yapan başlangıç koşulu. Sonuç: parçacık her turda aynı noktaya döner.
\item \emph{Başlangıç enerjisiyle ölç} ($\delta=0$): $f_0$ ve $\omega_0$.
\item \emph{Enerjiyi düzelt}: $\delta_1=-\omega_0/(\mathrm{d}\omega_h/\mathrm{d}\delta)$. Eğim
($-5.238\times10^6$~rad/s) kaçıklıktan bağımsız olduğu için önceden bilinir (soru 6). Bu $\delta_1$ ile
kapalı yörünge yeniden aranır ve $f_1$, $\omega_1$ ölçülür.
\item \emph{Ara değerle}: $f^*=f_0+(f_1-f_0)(0-\omega_0)/(\omega_1-\omega_0)$; yani $\omega_h=0$ olan
noktadaki $f$.
\end{enumerate}

\textbf{``Enerji korunumlu fırlatma''.} Yalnızca 2. adımdaki başlangıç tahminidir: momentum,
$\Phi=-E_0R_0\ln(R/R_0)$ potansiyeline göre fırlatma noktasındaki $x$'e uyarlanır. Soru 9'da görüleceği
gibi fırlatma noktası alansız bir bölgede olduğundan bu kuralın fiziksel anlamı yoktur. Önemli değildir:
prosedür $\delta$'yı ayrıca aradığı için sonuç bu başlangıç seçiminden bağımsızdır~\YOR. Ana rapordaki
açıklama bu yönde düzeltildi.

\textbf{İki nokta neden yeterli.} Çünkü $f$ ve $\omega_h$, ilgilenilen aralıkta $\delta$'ya doğrusaldır.
Beş noktalı tarama ($\delta=\pm10^{-6}$) bunu gösterdi~\BEL. İkinci noktada $|\omega_1|$, $|\omega_0|$'ın
yaklaşık $10^{-5}$'i kadardır (ilk tahmin zaten donmuş noktaya çok yakındır); yani ara değerleme çok
küçük bir aralıkta yapılır ve doğrusal olmamaktan gelen hata ihmal edilebilir~\BEL. Sınırlama: doğrusallık
yalnızca bir desende ($\sigma=10~\mu$m, seed~0) sınandı.

\soru{9}{Fırlatma noktası değişince sonucun değişmesini enerji farkıyla açıklayan argüman bence geçersiz.
Fırlatma noktası bir kuadrupolde. Bu argümanı yeniden tartışalım: bağımlılık gerçekten enerjiden mi geliyor?}

\textbf{İtiraz büyük ölçüde haklı; açıklama düzeltildi.} Kodda parçacık, bir hücrenin ilk sürüklenme
elemanının başından fırlatılır: ARC1 saptırıcısının hemen çıkışı, QF'ten önce (\kod{integrator.cpp},
eleman sırası 1, 2, \dots, 7, 0). Bu nokta alansız bir bölgededir. Ana rapordaki ``kapalı yörünge
dışarıdaysa potansiyel enerji farklıdır'' açıklaması orada geçersizdir: alansız bölgede potansiyel
$x$'e bağlı değildir.

\textbf{Doğru açıklama: sert kenarlı saptırıcılar.} Simülasyondaki saptırıcıların alanı, saptırıcı
sınırında aniden başlar ve sınırda potansiyel basamağı yoktur. Gerçek bir saptırıcıda saçak alanı vardır ve
toplam enerji $\gamma mc^2+q\Phi$ korunur. Sert kenarlı modelde ise parçacık saptırıcıya girerken potansiyel
basamağını tırmanmaz; saptırıcının içinde radyal hareketine göre iş yapılır ve bu iş girişte/çıkışta geri
alınmaz. Sonuç: \emph{aynı kapalı yörüngenin farklı hücrelerindeki alansız bölgelerde $|p|$ farklıdır}~\HIP.

\textbf{Doğrudan test~\BEL} (\kod{tests/teshis\_p\_azimut.py}, $\sigma=10~\mu$m, seed~0, CCW). Eski prosedürle
(hücre 0'da $|p|=p_{\rm sihirli}$) kapalı yörünge bulundu ve hücre 0, 6, 12, 18'de $|p|$ ölçüldü (ARC2
girişinde; ARC1 çıkışından ARC2 girişine kadar yalnız manyetik eleman olduğu için $|p|$, o hücrenin fırlatma
noktasındakiyle aynıdır):
\begin{center}\small
\begin{tabular}{lcccc}
\toprule
hücre & 0 & 6 & 12 & 18 \\
\midrule
$\sigma=10~\mu$m: $|p|/p_{\rm sihirli}-1$ & $+2\times10^{-12}$ & $-5.48\times10^{-7}$ & $-2.82\times10^{-7}$ & $-2.33\times10^{-7}$ \\
kaçıklık yok & $-1.6\times10^{-12}$ & $-1.6\times10^{-12}$ & $-1.0\times10^{-12}$ & $-1.6\times10^{-12}$ \\
\bottomrule
\end{tabular}
\end{center}
\sekil{ek_sekil_s4_firlatma.pdf}{Üst: bir FODO hücresinin eleman sırası ve fırlatma noktası. Sol alt: aynı kapalı yörüngenin farklı hücrelerinde $|p|$. Sağ alt: aynı desen için üç fırlatma kuralı; yalnız donmuş spin fırlatma noktasından bağımsız.}
Kaçıklık yokken fark sayısal sıfırdır; kaçıklıkla $\sim10^{-7}$'dir. Yani eski prosedür, fırlatma hücresi
değişince $\sim10^{-7}$ kadar farklı bir yörünge enerjisi seçiyordu. Soru 6'ya göre bu, yüzlerce kat
hedeflik bir fark yaratmaya yeter.

\textbf{Çapraz kontrol~\BEL.} Ana rapordaki dört donmuş spin koşumu (dört farklı fırlatma hücresinden)
kullanılarak donmuş spin yörüngesinin her hücredeki $|p|$'si hesaplandı. Hücre 0'a göre farklar
$-5.482$, $-2.816$, $-2.333\times10^{-7}$; doğrudan ölçümle \emph{dört basamak aynı}. Bu iki şey gösteriyor:
(1)~dört noktadan bulunan donmuş spin çözümleri tek ve aynı yörüngedir; (2)~yörünge boyunca $|p|$ değişimi
gerçektir ve hücreden hücreye fırlatma farkını tam olarak açıklar.

\textbf{Açık kalan nokta~\BEL.} Ana raporda eski sonuçların fırlatma noktasına göre değişimini
$\mathrm{d}f/\mathrm{d}\delta$ eğimiyle nicel olarak eşlemeye çalışmıştım; bu eşleme tutmuyor. Donmuş
noktanın dışında $f$'nin $\delta$'ya eğimi fırlatma noktasına göre farklı. Nedenini incelemedim; tahminim,
$\omega_h\ne0$ iken ölçülen $f$'ye karışan terimin, spinin başlangıç yönüne göre fırlatma noktasında
değiştiğidir~\HIP. Bu, donmuş spin sonuçlarını etkilemez.

\textbf{Gerçek makine için anlamı~\YOR.} Gerçek saptırıcıda saçak alanı enerjiyi korur ve alansız
bölgelerin hepsinde (aynı potansiyeldeyseler) $|p|$ aynıdır. Yani hücreden hücreye $|p|$ farkı sert kenar
modelinin bir artefaktıdır. Ama donmuş spin sonucu bu artefakttan etkilenmez, çünkü $|p|$'yi değil
$\omega_h=0$ koşulunu sabitler. Makaledeki sonuçlar için doğru tavır: sahte EDM'yi donmuş spin noktasında
vermek ve saçak alanının etkisini ayrıca sınamak (bu çalışmada yapılmadı).

\soru{10}{Eski prosedürle de $p\approx2$ bulunmuştu. O hâlde eski değerler nasıl 0.3--13 kat yanlış
olabilir? $\sigma^2$ ölçeklemesi neyi doğruluyor, neyi doğrulamıyor?}

Çünkü eski prosedürün hatası da $\sigma^2$ ile ölçeklenir. Adım adım~\BEL:
\begin{enumerate}[nosep]
\item Eski sonuç $=$ doğru sonuç $+$ (enerji kayması) $\times$ (eğim): $f_{\rm eski}=f^*+(\mathrm{d}f/\mathrm{d}\delta)\,\delta_{\rm fırlatma}$.
\item Enerji kayması kapalı yörüngeyle orantılıdır: $\delta_{\rm fırlatma}\propto\sigma$ (soru 9: kaçıklık
yokken sıfır).
\item Eğim de kaçıklıkla orantılıdır: $\mathrm{d}f/\mathrm{d}\delta\propto\sigma$ (soru 6).
\item Çarpımları $\propto\sigma^2$. Yani hata, doğru sonuçla aynı üsle ölçeklenir.
\end{enumerate}
Doğrudan görülüyor: enerji korunumlu başlangıç noktasındaki $f_0$'ın donmuş spin değerine oranı her seed için
$\sigma$'dan bağımsızdır: seed~0: $-4.57$, $-4.58$, $-4.59$; seed~1: $2.45$, $2.44$, $2.44$; seed~2:
$0.84$, $0.84$, $0.85$ ($\sigma=2.5$, 5, 10~$\mu$m)~\BEL. Oran 1'den farklı, ama sabit.

\sekil{ek_sekil_s5_sigma2.pdf}{Sol: donmuş spin ve sabit fırlatma kuralı ile $|f|$; ikisi de $\sigma^2$ ile ölçeklenir. Sağ: oranları her seed için sabittir ama 1'den farklıdır; $\sigma^2$ testi bu yanlılığı yakalayamaz.}
Yani $\sigma^2$ testi \emph{mekanizmanın türünü} doğrular (ikinci mertebe, geometrik faz, birinci mertebe
sızıntı yok); \emph{mutlak değeri} doğrulamaz. Mutlak değer için ek bir değişmezlik testi gerekir (fırlatma
noktasından bağımsızlık gibi)~\YOR.

%======================================================================
\section{Şekil 4 ve TSVD eşiği (ana rapor 3.5)}
%======================================================================

\soru{11}{Sıfır BPM ofseti gerçekçi mi? Ofset olunca bu tablo nasıl değişir?}

Mutlak anlamda gerçekçi değildir: BPM ofsetleri tipik olarak yüzlerce mikrometredir ve BBA olmadan
bilinmez. Ama Şekil~4 doğru yorumlanırsa gereksiz bir varsayım değildir: yöntem BBA referansına göre
\emph{fark} ölçer (soru 5); sabit ofset farkta sadeleşir. Bu yorumda ``sıfır ofset'' yerine doğru
varsayım ``referanstan beri ofset değişmedi''dir. Ofsetin değişimi $\delta$ olarak T2 ve Adım~6'da
modellendi: rijit BPM ve iki demetle $\sigma_\delta\le1~\mu$m'de sahte EDM sayısal tabanda kalıyor~\BEL.
Lab-sabit BPM'de ofset değişiminin etkisi ayrı olarak ölçülmedi.

Mutlak ofsetle (100~$\mu$m, referans yok) aynı tablo medyan 198 kat kötüleşir (soru 5). Makalede Şekil~4
``BBA'dan sonra, referansa göre'' koşulunda verilmeli ve bu açıkça yazılmalıdır~\YOR.

\soru{12}{\kod{rcond}~$=0.01$ eşik seçimi sonucu ne kadar belirliyor?}

Çok. Ana rapor Adım~1'den (5~$\mu$m simetrik, 14~nm, 200 seed, kalan simetrik rms)~\BEL:
\begin{center}\small
\begin{tabular}{lccc}
\toprule
durum & \kod{rcond}~$=0.01$ & en iyi eşik & en iyi eşik değeri \\
\midrule
$\rho=0$, tek demet (E1) & 3.87~$\mu$m & 0.051~$\mu$m & 0.003 \\
$\rho=0$, iki demet (E3) & 0.52~$\mu$m & 0.017~$\mu$m & 0.003 \\
$\rho=1$, iki demet, $\rho$ biliniyor (E4) & 1.30~$\mu$m & 0.012~$\mu$m & 0.0003 \\
\bottomrule
\end{tabular}
\end{center}
0.01 ile tek demet hiç çalışmıyor gibi görünür (\%23 düzeltme); 0.003 ile 100 kat düzeltiyor. Şekil~4'teki
``tek demet 12, iki demet 0.0015'' farkı (yeni prosedür) büyük ölçüde bu eşikten gelir: Adım~1'de uygun
eşikle tek demet de hedefin $\sim$120 kat altına iniyor ($C=0.008$)~\BEL. Makalenin ``tek demet kördür''
iddiası bu eşikle kurulmuş bir karşılaştırmaya dayanıyor~\YOR.

%======================================================================
\section{Adım 1: BPM modeli (ana rapor 4)}
%======================================================================

\soru{13}{``Simetrik kayma, tek demette bile büyük bir sinyal olarak görünür.'' Bu sezgiyi açar mısın?
Makalenin iki demet gerekçesini zayıflatıyor mu?}

\sekil{ek_sekil_s6_bpm.pdf}{BPM modeli. (a) Lab-sabit BPM demetin laboratuvardaki konumunu, (b) rijit BPM demetin kuadrupol merkezine göre konumunu okur.}
Lab-sabit BPM demetin \emph{laboratuvardaki} konumunu okur: $r=x$. Simetrik kaçıklığın ürettiği yörünge
küçüktür (soru 1), bu yüzden okuma küçüktür. BPM kuadrupole rijit bağlıysa ve kuadrupolle birlikte $d$ kadar
kayarsa, demetin \emph{BPM'e göre} konumunu okur: $r=x-d$. Demet pek kımıldamadığı için ($x$ küçük)
$r\approx-d$ olur; yani kaçıklığın kendisini doğrudan okur. Matris dilinde okuma $(R-I)d$'dir ve $-I$ terimi
her kaçıklık desenini, simetrik dahil, büyük bir sinyale çevirir.

Makale açısından anlamı~\YOR: ``Simetrik bileşen yörüngeye kördür, iki demet gerekir'' gerekçesi yalnızca
lab-sabit BPM'de geçerlidir. Rijit BPM'de tek demet de görür (Adım~1: E2, $C=0.006$). Ama iki demetin
\emph{üstünlüğü} kaybolmaz, başka bir yerden gelir: BPM--kuadrupol göreli kayması $\delta$ varken iki
demet çok daha sağlamdır (T2: $\sigma_\delta=0.3$ ve 1~$\mu$m'de 24 ve $\sim$800 kat). Yani gerekçe
``görünürlük''ten ``sağlamlık''a kayar. Bu, makaleyi zayıflatmaz ama hikâyesini değiştirir (soru 41).

\soru{14}{``En iyi eşik'' neye göre seçildi? Gerçek makinede doğru kaçıklık bilinmeyeceği için bu eşik
pratikte nasıl seçilir? Bu karşılaştırma adil mi?}

\textbf{Nasıl seçildi.} Bir eşik ızgarası ($10^{-1}$--$10^{-5}$) tarandı; her kestirici için 200 seed'in
medyan kalan simetrik rms'ini en küçük yapan eşik alındı. Bu, \emph{gerçek kaçıklığı bilen} bir seçimdir
(``kahin'' seçim). Her kestiriciyi kendi en iyi koşulunda karşılaştırdığı için kestiriciler arasında adildir;
ama gerçek makinede yapılamaz.

\textbf{Pratikte: uyuşmazlık ilkesi (Morozov).} Gerçek kaçıklık gerekmez; yalnız BPM gürültüsü $\sigma_n$
bilinmelidir (tekrarlı ölçümlerden kolayca ölçülür). Kural: fit artığını gürültü düzeyine indiren en büyük
eşiği seç, yani $\|S\hat d-y\|\le\tau\,\sigma_n\sqrt{N}$ ($\tau=1.1$, $N$ ölçüm sayısı). Her seed kendi
eşiğini kendi okumasından seçer. Test (\kod{ek\_rcond\_secimi.py}, 200 seed, kalan simetrik rms)~\BEL:
\begin{center}\small
\begin{tabular}{lccc}
\toprule
durum & uyuşmazlık ilkesi & kahin & en sık seçilen / kahin eşik \\
\midrule
$\rho=0$, E1 & 0.052~$\mu$m & 0.051~$\mu$m & 0.003 / 0.003 \\
$\rho=0$, E3 & 0.017~$\mu$m & 0.017~$\mu$m & 0.003 / 0.003 \\
$\rho=1$, E2 & 0.20~$\mu$m & 0.115~$\mu$m & 0.0003 / 0.0003 \\
$\rho=1$, E4 & 0.017~$\mu$m & 0.012~$\mu$m & 0.003 / 0.0003 \\
$\rho=1$, E3 (yanlış model) & 6.6~$\mu$m & 4.4~$\mu$m & $10^{-5}$ / 0.03 \\
\bottomrule
\end{tabular}
\end{center}
\sekil{ek_sekil_s7_rcond.pdf}{Üç eşik seçimi: sabit 0.01, uyuşmazlık ilkesi (gerçekte uygulanabilir) ve kahin. Model doğruyken uyuşmazlık ilkesi kahine çok yakın.}

Model doğruyken uyuşmazlık ilkesi kahine çok yakın (en kötü 1.7 kat). Ana rapordaki karşılaştırmalar
bu yüzden gerçekçidir~\YOR. Model yanlışken (son satır) ilke hiçbir eşikte gürültü düzeyine inemez ve en
küçük eşiğe gider. Bu bir kusur değil, \emph{yararlı bir uyarıdır}: ``fit gürültü düzeyine inmiyor'',
modelin (ör.\ BPM modelinin) yanlış olduğunu veriden gösterir~\YOR.

\soru{15}{$\rho$'nun bilindiği varsayılıyor. Bu bilgi gerçekte nereden gelir? Ölçülebilir mi? $\rho$
yanlış bilinirse ne olur?}

\textbf{Nereden gelir~\YOR.} İki yol var. (a)~Mekanik tasarım: BPM kuadrupolün içine ya da aynı rijit
desteğe monte edilmişse $\rho\approx1$ tasarımdan bilinir. (b)~Ölçüm: kuadrupol hareket ettiricilerle
bilinen bir miktar kaydırılır (ya da yer hareketi sensörleriyle izlenir) ve BPM okumasının değişimi
modelin öngördüğüyle karşılaştırılır. $\rho=1$ için okuma değişimi $(R-I)\Delta d$, $\rho=0$ için
$R\,\Delta d$'dir; fark büyük ($-I$ terimi) olduğundan $\rho$ birkaç yüzde hassasiyetle ölçülebilir~\HIP.
Ayrıca soru 14'teki uyuşmazlık ilkesi yanlış $\rho$'yu kendiliğinden işaret eder.

\textbf{Yanlış bilinirse~\BEL.} Gerçek $\rho=1$, kestirici $\hat\rho\ne1$ kullanıyor (5~$\mu$m simetrik,
14~nm; doğrusal: 200 seed; sahte EDM: 6 seed):
\begin{center}\footnotesize\setlength{\tabcolsep}{3pt}
\begin{tabular}{lcccccc}
\toprule
& \multicolumn{3}{c}{iki demet (E4)} & \multicolumn{3}{c}{tek demet (E2)} \\
\cmidrule(lr){2-4}\cmidrule(lr){5-7}
$\hat\rho$ & sim [$\mu$m] & $m$=2--3 & $C$ & sim [$\mu$m] & $m$=2--3 & $C$ \\
\midrule
0.80 & 1.95 & 0.008 & -- & 2.18 & 0.28 & -- \\
0.90 & 1.94 & 0.004 & \Ceqnine & 0.84 & 0.19 & \Cetnine \\
0.95 & 0.60 & 0.006 & \Ceqnf & 0.40 & 0.14 & \Cetnf \\
0.98 & 0.20 & 0.007 & -- & 0.20 & 0.09 & -- \\
1.00 & 0.013 & 0.008 & 0.003 & 0.12 & 0.05 & 0.006 \\
\bottomrule
\end{tabular}
\end{center}
\sekil{ek_sekil_s8_rhohat.pdf}{Gerçek $\rho=1$ iken kestiricinin yanlış $\hat\rho$ varsayması. Sol: kalan simetrik rms (doğrusal, 200 seed, gölge 16--84\%). Sağ: C++ sahte EDM (nokta = seed).}
($\hat\rho=1$ satırındaki $C$ değerleri ana rapor Adım~1'den, 10 seed medyanı.) \RhoYorum

\soru{16}{Eski ve yeni sonuç arasındaki bu kadar büyük fark neden? $\rho=0.5$'in ölçülmemesi bir eksik mi?
Ölçülürse ne bekliyoruz?}

\textbf{Fark neden.} $\rho=1$'de E3'ün bıraktığı artık büyüktür (4.4~$\mu$m) ama ağırlıklı olarak spinin az
duyarlı olduğu harmoniklerdedir (soru 17). Bu tür bir artık, kapalı yörüngeyi yine de belirgin biçimde
bozar; eski prosedürde bu, seçilen yörünge enerjisini kapalı yörüngeyle orantılı biçimde kaydırır (soru 9, 10). Sonuç:
gerçek değeri $\sim0.2$ olan bir durum, enerji yanlılığıyla 0.5--3 olarak ölçülmüştü~\YOR.

\textbf{$\rho=0.5$, şimdi ölçüldü~\BEL} (6 seed, her kestiricinin en iyi eşiği; doğrusal kalan simetrik rms
200 seed medyanı):
\begin{center}\small
\begin{tabular}{lccc}
\toprule
kestirici & kalan sim [$\mu$m] & $C$ medyan [aralık] & $C>1$ \\
\midrule
E1 (tek demet, BPM sabit sanılır) & 4.50 & \RhoBirEbir \\
E2 (tek demet, $\rho$ biliniyor) & 0.29 & \RhoBirEiki \\
E3 (makale, BPM sabit sanılır) & 3.30 & \RhoBirEuc \\
E4 (iki demet, $\rho$ biliniyor) & 0.009 & \RhoBirEdort \\
\bottomrule
\end{tabular}
\end{center}
\RhoBesYorum

%======================================================================
\section{Adım 1 yorumu (ana rapor 4.5)}
%======================================================================

\soru{17}{4.5 bölümünü daha ayrıntılı açıklar mısın? E3 neden $\rho=1$'de simetrik sürüklenmenin çoğunu geri
bulamadığı halde sahte EDM'yi hedefin altına indiriyor? ``33 kat ile 4 kat arasındaki fark'' nereden
gelebilir?}

\textbf{Durum.} Gerçekte BPM'ler kuadrupolle birlikte hareket ediyor ($\rho=1$, okuma $(R-I)d$); E3 ise
BPM'leri sabit sanıyor (model $R\,d$). Kestirim $\hat d=S^{+}[(R_{\rm CCW}-I);(R_{\rm CW}-I)]\,d$ olur; yani
yanlış bir doğrusal harita.

\textbf{Neden yine de iyi.} Bu yanlış harita her harmoniği farklı etkiler~\BEL\ (ana rapor 4.5, kalan güç /
başlangıç gücü): $m=2$--3'te 0.22--0.24, $m=0$--1'de 1.1--1.4, $m\ge9$'da $\approx1$. Yani E3, yörünge
tepkisinin büyük olduğu harmonikleri (betatron tune'u $Q\approx2.7$ yakınında) düzeltiyor, diğerlerine
dokunmuyor. Nedeni~\YOR: $m=2$--3'te yörünge tepkisi $R$ büyüktür, $R-I\approx R$ olur, yanlış model
neredeyse doğrudur. Uzak harmoniklerde $R$ küçüktür ve $-I$ baskındır; orada yanlış model kestirimi
bozar. Spin ise tam olarak $m=2$--3'e duyarlıdır (Adım~5; demet $m=0$'da kuadrupolleri izler ve
merkezlerinden geçer). Bu yüzden 4.4~$\mu$m'lik artık büyük ama zararsıza yakındır.

\sekil{ek_sekil_s9_e3_harmonik.pdf}{$\rho=1$'de makale kestiricisinin (E3) her harmonikteki etkisi: kalan güç / başlangıç gücü. Yalnız spin-duyarlı $m=2$--3 belirgin biçimde azalıyor.}
\textbf{33 kat ile 4 kat arasındaki fark (açık; tahminler).} Sahte EDM bilineerdir: $f\approx\sum_{ij}
x_i\,K_{ij}\,y_j$ (yatay ve düşey kaçıklık desenlerinin çarpımı). Güç oranı her düzlemi ayrı ayrı ölçer;
$C$ ise iki düzlemin \emph{çarpımına} bağlıdır. İki olası katkı~\HIP:
\begin{enumerate}[nosep]
\item Her iki düzlemde genlik $\sim\sqrt{0.23}\approx0.48$ kat azalırsa çarpım $\approx0.23$ kat azalır;
bu tek başına 4 kat verir. Ancak $C$, $m=2$--3 içinde de belirli faz ilişkilerine (yatay ve düşey desenlerin
aynı/dik fazlı bileşenlerine) farklı duyarlıdır (Adım~5'teki ``aynı faz'' ve ``dik faz'' sütunları 3--10
kat farklı). E3 bu bileşenlerden spin açısından en etkili olanı daha çok bastırıyorsa fark açıklanır.
\item Kalan artıkta $m=0$--1 ile $m=2$--3 arasındaki çapraz terimler başlangıçtakinden farklı işaretli
olabilir ve kısmen sadeleşebilir.
\end{enumerate}
Bunları ayırmak için bilineer çekirdek $K$'nın harmonik tabanda ölçülmesi gerekir (Adım~5'in genişletilmiş
hali; yapılmadı).

%======================================================================
\section{Adım 2: Saptırıcı hataları (ana rapor 5)}
%======================================================================

\soru{18}{Düşey eksen burada $z$ olarak yazılmış; projede düşey eksen $y$. Bu $E_y$ olmalı. Roll'un öncü
teriminin neden $E_0\phi$ olduğunu da açıklar mısın?}

Haklısın; ana raporda $E_y\approx E_0\phi$ olarak düzeltildi.

Neden $E_0\phi$: Saptırıcının alanı referans yörüngede radyaldir, $\vec E=E_0\,\hat x$. Saptırıcı, demet
ekseni ($s$ ekseni) etrafında küçük bir $\phi$ açısıyla dönerse alan deseni de aynı açıyla döner:
$\vec E=E_0(\cos\phi\,\hat x+\sin\phi\,\hat y)$. Küçük açıda $E_x\approx E_0$ (değişim $\phi^2$ mertebesi)
ve $E_y\approx E_0\phi$. Yani roll, birinci mertebede, yörünge boyunca düzgün bir düşey elektrik alanı
ekler. Bu alan demete düşey bir tekme verir: $\theta_E=(L/R_0)\,\phi=(2\pi/48)\,\phi$ ($L$ saptırıcı boyu,
$L/R_0$ saptırıcının bükme açısı)~\YOR.

\sekil[0.45]{ek_sekil_s10_roll.pdf}{Roll: alan vektörü $\phi$ kadar döner; düşey bileşen birinci mertebede $E_0\phi$.}
\soru{19}{Asal açı kosinüsü ne demek ve 0.99'un üstünde olması pratikte neyi gösteriyor? Bu yakın
dejenerelik fiziksel olarak nereden geliyor?}

\textbf{Tanım.} İki alt uzay düşün: A, bütün olası saptırıcı hatalarının ürettiği iki demetli okumalar
($96$ boyutlu uzayda 48 boyutlu bir alt uzay); B, bütün simetrik kuadrupol kaçıklıklarının ürettiği
okumalar. Asal açılar, bu iki alt uzay arasındaki en küçük açılardır: önce A'dan ve B'den birer birim
vektör seçilir, aralarındaki açı en küçük olacak biçimde; sonra bunlara dik olanlar arasında aynı şey
tekrarlanır. Kosinüs 1 ise o doğrultuda iki alt uzay çakışır.

\textbf{0.99'un anlamı.} Kosinüs 0.99, açı $\approx8^\circ$ demektir. Pratikte: bazı saptırıcı hata
desenleri, bazı simetrik kuadrupol desenleriyle iki demette de neredeyse aynı okumayı üretir. Gürültü
olmasa bile bunları ayırmak için okumadaki \%1--14'lük farka güvenmek gerekir; gürültü bu farkı
büyütür. Bu yüzden saptırıcı sütunlu fitin koşul sayısı 73--118'den $5.6$--$7.9\times10^3$'e çıktı~\BEL.

\textbf{Fiziksel kaynağı~\YOR.} Halkada kapalı yörünge tepkisi, kaynağı ne olursa olsun, betatron tune'una
yakın harmoniklerde ($m\approx2$--3) rezonansla büyür ($G_k\propto1/|Q^2-k^2|$). Hangi elemandan gelirse
gelsin, dağılmış bir tekme deseninin yörüngesi bu birkaç şekil tarafından domine edilir. Bu yüzden saptırıcı
ve kuadrupol desenleri bu birkaç doğrultuda aynı yörünge şeklini üretir. Kanıt: yakın-çakışan doğrultular
çiftler halinde gelir (0.9987, 0.9986; 0.9965, 0.9964) --- aynı harmoniğin cos ve sin fazları --- ve aynı
yakın dejenerelik \emph{antisimetrik} alt uzayla da vardır (x'te 4 doğrultu)~\BEL. Yani mesele simetriğe
özgü değil, rezonant yörünge tepkisine özgüdür.

\soru{20}{Roll etkisi ($10^8\times$hedef) fiziksel olarak ne anlama geliyor? CW$-$CCW farkında $10^{-10}$
hassasiyetle sadeleşmesine güvenebilir miyiz? CW/CCW simetrisi bozulursa ne kadar tehlikeli?}

\textbf{Anlamı.} $3\times10^{-7}$ rms roll için $f\approx0.13$~rad/s: spin $\sim$8 saniyede 1 radyan
düşeye döner. Mekanizma (ana rapor 5.6): roll'un düşey elektrik tekmelerini kapalı yörünge kuadrupollerin
manyetik tekmeleriyle dengelemek zorundadır ve manyetik alan spini momentumdan $G\gamma$ kat fazla
döndürür. Bu, EDM deneylerinde bilinen ``ortalama radyal manyetik alan'' sistematiğinin elektrik
roll'dan doğan biçimidir~\YOR.

\textbf{Sadeleşme neden olur.} Elektrik tekmeleri iki yönde aynıdır; kuadrupol tekmeleri de, aynı kapalı
yörüngeyi dengelemek için, iki yönde aynı spin etkisini verir. Sonuç $f_{\rm CW}=f_{\rm CCW}$ (test paketi:
9 basamak) ve gerçek EDM ise iki yönde zıt işaretlidir. Simülasyonda iki yön tam simetrik olduğu için
$10^{-10}$ bağıl sadeleşme görülür. Bu, sayısal hassasiyetin göstergesidir; gerçek makinenin değil~\YOR.

\textbf{Simetri bozulursa~\HIP.} Kalan $C\approx\epsilon_a\times|f_{\rm roll}|$, burada $\epsilon_a$
iki yönün bağıl asimetrisidir. $C<1$ için $\epsilon_a<1/(1.26\times10^8)\approx8\times10^{-9}$ gerekir
($3\times10^{-7}$ roll için). Asimetri kaynakları: iki demetin farklı enerjide donması (her biri kendi
$\omega_h=0$ noktasında), CW/CCW için farklı kapalı yörüngeler (manyetik elemanlar yönle işaret değiştirir),
farklı demet yoğunlukları. Örneğin $f_{\rm roll}\propto G\gamma$ ve iki yönün enerji farkı $\Delta\delta$
ise $\epsilon_a\sim\Delta\delta$ mertebesindedir; $\Delta\delta=10^{-8}$ bile sınırdadır. Statik roll
(0.1~mrad, Hakem A) $3\times10^{-7}$'den 300 kat büyük olduğundan $f_{\rm roll}$ da 300 kat büyük olur;
bu, statik roll'un $\mu$rad düzeyinde hizalanması ya da düşey elektrik trimiyle telafi edilmesi gerektiğini
gösterir (soru 45). Bu, ana raporda ayrıca test edilmedi; yöntemin değil deneyin genel bir gereksinimidir
ve Omarov vd.\ CW/CCW + polarite çevirme tekniklerini bunun için önerir~\cite{omarov}.

%======================================================================
\section{Adım 3: Düzelticiler (ana rapor 6)}
%======================================================================

\soru{21}{Diğer düzelticiler neden çalışmıyor? Fiziksel nedeni?}

Amaç kaçıklığı ``iptal etmek''tir: demetin her kuadrupolde gördüğü alan, hizalı bir kuadrupoldekiyle aynı
olmalı. Üç düzeltici türüne bu açıdan bakalım~\YOR:
\begin{itemize}
\item \textbf{Mq (kuadrupol içi manyetik).} Alanı, kaçıklığın yarattığı dipol alanıyla \emph{aynı yerde,
aynı türde} siler: $g(y-d)+g\hat d\approx g\,y$. İki demet için de, spin için de kuadrupol hizalı gibidir.
\item \textbf{Md (saptırıcıda manyetik).} Tekmesi yönle işaret değiştirir, yani yalnız \emph{fark kanalına}
etki eder. Simetrik sürüklenme ise iki demette aynı işaretli bir düzeltme ister (demeti iki yönde de aynı
tarafa, kuadrupol merkezine doğru kaydırmak). Md bunu yapamaz. ``Sıfır'' hedefiyle BPM okumasını sıfırlar
ama demet kuadrupol merkezlerinden geçmez; kuadrupollerdeki dipol alanları kalır ve sahte EDM sürer.
``Merkez'' hedefi ise iki yönde zıt yönlü bir sapma ister ki bu da kendi sahte EDM'sini üretir.
\item \textbf{Ed (saptırıcıda elektrostatik).} Tekmesi iki yönde aynıdır, yani toplam kanalına etki eder;
demeti iki yönde de kuadrupol merkezlerinden geçirebilir (ana rapor: 0.07--0.12~$\mu$m). Ama bunun için
büyük elektrik tekmeleri gerekir ($\varepsilon$ ve $\phi$ düzeltmeleri $\sim7.5\times10^{-6}$ rms~\BEL).
Bu tekmeler yeni bir kapalı yörünge şekli dayatır ve roll etkisine benzeyen birinci mertebe bir spin
dönmesi (çift kısım) ile kendi $x_{\rm CO}y_{\rm CO}$ geometrik fazını (C) üretir (soru 22).
\end{itemize}
Özetle: sorun kuadrupolün kendi alanındadır; onu yalnız kuadrupolün yerinde, manyetik bir alan tam
iptal edebilir.

\soru{22}{Elektrostatik düzeltici bu mekanizmayla nicel olarak ne kadar sahte EDM üretir? Tablo 6'daki Ed
değerleri bununla açıklanıyor mu? Yeni prosedürle ölçüldü mü?}

Yeni prosedürle şimdi ölçüldü~\BEL\ (Adım~3 senaryosu: 5~$\mu$m simetrik, lab-sabit BPM, 14~nm, iki demetli
kestirim; 5 seed):
\begin{center}\footnotesize\setlength{\tabcolsep}{3pt}
\begin{tabular}{lccccc}
\toprule
durum & düzeltici ($\varepsilon$ / $\phi$) & demet$-$merkez & $C$ medyan [aralık] & $C>1$ & $|$çift$|$ \\
\midrule
Ed, sıfır & \EdSifirKuv & 4.8~$\mu$m & \EdSifirC \\
Ed, merkez & \EdMerkezKuv & 0.07--0.12~$\mu$m & \EdMerkezC \\
\midrule
eski (Tablo 6), sıfır & & & 14, 4.4, 6.7 & & $2\times10^7$--$1.5\times10^8$ \\
eski (Tablo 6), merkez & & & 2.9, 1.5, 2.8 & & $10^6$--$9\times10^6$ \\
\bottomrule
\end{tabular}
\end{center}
\EdYorum

%======================================================================
\section{Adım 4 ve tartışma testleri (ana rapor 7--8)}
%======================================================================

\soru{23}{Hakem A'nın endişesi tam olarak neydi ve bu sonuç ona nasıl karşılık geliyor? Yakalanamayan
kısmın spine etkisinin neden daha küçük olduğunu açıklar mısın?}

\textbf{Endişe.} Hakem A iki şey söylüyor: (1)~gerçek hizalama hataları korelasyonludur; iki kuadrupol
arasındaki göreli hata aradaki mesafenin kareköküyle büyür (yer hareketi, ATL yasası); (2)~uzun dalgalı
kaymada demet kuadrupolleri izler, yörünge bastırılır, ama spinin aynı bastırmayı görüp görmediği belli
değildir.

\textbf{Karşılık~\YOR.} Uzun dalgalı yer hareketi kuadrupolleri \emph{ve} yanlarındaki BPM'leri birlikte
taşır. Bu, $\rho\approx1$ demektir. Makalenin modeli (lab-sabit BPM) bu durumda simetrik sürüklenmenin
yalnız \%30'unu yakalar (ana rapor Tablo 8)~\BEL. Hakemin ``uzun dalgalı hareket'' kaygısının asıl
karşılığı budur: sorun yörüngenin bastırılması değil, BPM'lerin de hareket etmesidir.

\textbf{Spine etkisi neden küçük.} Yakalanamayan kısım, uzun dalgalı (düşük $m$) kısımdır ve demetin
kuadrupolleri izlediği harmoniklerdir. Demet kuadrupolleri izliyorsa onların merkezlerinden geçer, alan
görmez, spin dönmez. Ölçülen: $m=0$'da birim genlik başına sahte EDM, $m=2$'dekinden $\sim10^4$--$10^5$
kat küçük (ana rapor Adım~5, eski prosedür)~\BEL. Doğrudan bir kanıt da var: $\rho=1$'de E3 4.4~$\mu$m
artık bıraktığı halde $C=0.19$ veriyor (soru 17)~\BEL. Yani hakemin ikinci sorusunun cevabı ``evet, spin de
bastırmayı görür; ama yalnız gerçekten izlenen harmoniklerde'' olur.

\soru{24}{Adım 1--4 sonunda iki sorunun neden kaldığını ve diğerlerinin neden çözülmüş sayıldığını
özetler misin?}

\begin{center}\small
\begin{tabularx}{\textwidth}{>{\raggedright\arraybackslash}p{3.6cm}>{\raggedright\arraybackslash}p{1.6cm}>{\raggedright\arraybackslash}X}
\toprule
konu & durum & neden \\
\midrule
düzeltici türü (Adım 3) & çözüldü & tasarım seçimiyle: Mq kullanılırsa sorun yok \\
korelasyonlu hata (Adım 4) & çözüldü & doğrusal kestirici desen modelinden bağımsız; korelasyon yalnız simetrik payı artırıyor \\
model hatası, $\beta$-beat (Adım 1) & kısmen & iki demet üstünlüğünü azaltır, ama sahte EDM hedefin altında kalıyor \\
\textbf{BPM hareketi} (Adım 1, 4) & kalmıştı & makale modeli bunu hiç içermiyordu; çözüm modeli değiştirmeyi (rijit BPM + $\rho$-bilen kestirici) gerektirdi \\
\textbf{saptırıcı sürüklenmesi} (Adım 2) & kalmıştı & hata toplam kanalına, yani simetrik sinyalin kanalına düşüyor; ölçüme yeni bilinmeyenler eklemek ve donanım (trim) gerektirdi \\
\bottomrule
\end{tabularx}
\end{center}
İki sorunun ortak özelliği, ikisinin de yöntemin \emph{ölçüm modelinin} dışında kalmasıdır: biri ölçüm aletinin
(BPM) kendisinin hareketi, diğeri simetrik sinyalle aynı kanala düşen ikinci bir hata kaynağı. Diğerleri
ya bir tasarım seçimiyle ya da kestiricinin doğrusallığıyla kendiliğinden çözülüyordu~\YOR.

\soru{25}{Saptırıcı trimleri donanım olarak nasıl yapılır, gerçekçi mi? Gerekli gerilim çözünürlüğü ve
kararlılığı ne olur?}

\textbf{Nasıl~\YOR.} (a)~Alan trimi ($\varepsilon$): her saptırıcı plakasına ana yüksek gerilim kaynağıyla
seri bağlı küçük, yüzer bir trim kaynağı; ya da her saptırıcı için ayrı ana kaynak. (b)~Roll trimi
($\phi$, düşey alan): üst ve alt plaka (ya da plaka yarıları) arasında küçük bir gerilim farkı. Bu, düşey
bir alan $E_y=\Delta V/g$ üretir ($g$ plaka aralığı).

\textbf{Gereken büyüklük~\HIP.} Senaryodaki sürüklenme $\varepsilon,\phi\sim3\times10^{-7}$ rms. Alan trimi
için bu, ana gerilimin $3\times10^{-7}$'si demektir; $\pm200$~kV'de $\sim0.06$~V. Düşey trim için
$E_y=E_0\phi\approx4.4~\text{MV/m}\times3\times10^{-7}\approx1.3$~V/m; 3~cm aralıkla $\Delta V\approx0.04$~V
(plaka aralığı varsayım). Trimin kendisi küçük ve kolaydır.

\textbf{Asıl sorun: ana gerilimin kararlılığı~\LIT.} COSY'deki 200~kV kaynakların dalgalanması (ripple)
tepeden tepeye $\sim10^{-5}$'tir ve EDM prototip halkası için daha iyisinin geliştirilmesi gerektiği
yazılıdır~\cite{cpedm_feas}. nEDM deneyi için yapılan iki kutuplu 200~kV kaynak \%0.01 (100~ppm) kararlılık
ve 5~ppm dalgalanma belirtir~\cite{fug}. Başka bir halkadaki yüksek gerilim sisteminde 10~ppm kararlılık
raporlanmıştır~\cite{hirfl}. Yani bugünkü teknoloji $10^{-5}$--$10^{-4}$ düzeyindedir; bizim senaryomuzdaki
$3\times10^{-7}$ sürüklenme bundan 30--300 kat küçüktür.

\sekil{ek_sekil_s11_hv.pdf}{Saptırıcı alan kararlılığı: bu çalışmanın varsayımı ve bulunan literatür değerleri (logaritmik eksen).}
\textbf{Bu ne kadar ciddi~\YOR.} Birkaç ayrım gerekir:
\begin{itemize}[nosep]
\item Bütün saptırıcılarda \emph{ortak} değişim enerjiyi kaydırır; donmuş spin geri beslemesi bunu
otomatik düzeltir. Önemli olan saptırıcıdan saptırıcıya \emph{bağımsız} değişimdir.
\item Hızlı dalgalanma (kHz) yörüngeyi titretir ama spin ve kestirim zaman ortalaması alır; önemli olan
yavaş sürüklenmedir.
\item Bağımsız yavaş sürüklenme $10^{-5}$ olsaydı, Adım~2'nin doğrusallığına göre her saptırıcı
$\sim40~\mu$m yörünge bozulması yaratırdı ve T1'deki $\varepsilon\times\phi$ etkisi (amplitüdün karesiyle)
$\sim(33)^2\approx1000$ kat büyürdü. Trim sonrası $C\approx0.01$ olan sonuç $\sim10$'a çıkardı.
\end{itemize}
Bu yüzden saptırıcı başına alan kararlılığı, yöntemin (aslında herhangi bir proton EDM halkasının)
\textbf{en zor tasarım gereksinimidir} (soru 40). Bu ana raporda yeterince vurgulanmamıştı.

\soru{26}{Kestirim bu kadar kötüyken yöntem neden yine de çalışıyor? ``Kötü kestirilen doğrultular
yörüngeye zaten küçük katkı verir'' argümanını açar mısın?}

Kestirim problemi $y=D\,e+n$ ($D$ saptırıcı tepki matrisi, $e$ saptırıcı hataları). $D$'nin tekil değer
ayrışımında büyük tekil değerli doğrultular, yörüngeyi güçlü etkileyen hata desenleridir; bunlar iyi
kestirilir. Küçük tekil değerli doğrultular yörüngeye neredeyse hiç iz bırakmaz; gürültü altında kalır,
TSVD onları atar ve kestirim orada sıfır olur. Tek tek saptırıcı hataları her iki tür doğrultudan da pay
içerir; bu yüzden ``saptırıcı başına'' hata trimden sonra \%25--80 kalır~\BEL.

Ama sahte EDM, saptırıcı hatalarının \emph{ürettiği yörünge} üzerinden gelir ($\varepsilon\times\phi$
etkisi $\approx x_{\rm CO}\,y_{\rm CO}$). Kalan hata yalnız yörüngeye zayıf etki eden doğrultulardadır; yani
kalan yörünge küçüktür. Ölçülen: saptırıcı kaynaklı yörünge \%97--99 azalıyor~\BEL. $C$ yörüngelerin
çarpımıyla orantılı olduğundan, her düzlemde \%98 azalma çarpımda $\sim2500$ kat azalma demektir~\YOR.
Kısacası: ölçemediğimiz şey, zaten zarar veremeyen şeydir.

\soru{27}{$\delta$'nın bileşenleri (mekanik kayma ve BPM elektroniğinin ofset kayması) ayrı ayrı ne
büyüklükte beklenir? Hangisi baskın olur?}

\textbf{Elektronik~\LIT.} SLS~2.0 BPM sistemi için elektronik sürüklenme şartı 0.1~$\mu$m/saat
belirtilmiştir~\cite{sls2}. APS'te kısa sürede 2~$\mu$m/saat sürüklenme ölçülmüş ve $\sim$1.5 saatte
oturmuştur~\cite{aps_drift}. SPring-8 yükseltmesi için yeni elektronik prototipi uzun süreli kararlılıkta
1~$\mu$m içinde kalmıştır; aynı çalışma BPM için 5~$\mu$m uzun süreli kararlılık şartı koyar~\cite{spring8}.
APS MBA yükseltmesi, 7 günlük zaman ölçeğinde $\le1~\mu$m sürüklenme ister~\cite{aps_mba}.

\textbf{Mekanik~\HIP.} Çeliğin ısıl genleşmesi $\sim12\times10^{-6}$/K'dir. BPM elektriksel merkezi ile
kuadrupol manyetik merkezi arasında $\sim10$~cm'lik bir kol varsa, iki parçanın sıcaklık farkındaki 0.1~K'lik
değişim $\sim0.1~\mu$m kayma verir. Bir LEP~2 çalışmasında, gerili tel sensörleriyle izlenen BPM bloklarında
dolum sırasında birkaç mikronluk görünür kaymalar raporlanmıştır~\LIT~\cite{lep}. SRRC'de BPM'in girder'a
göre konumunun bir vardiyada 1~$\mu$m içinde kaldığı ve sonra 0.1~$\mu$m'ye iyileştirildiği raporlanır~\cite{lep}.

\textbf{Hangisi baskın~\YOR.} Sıcaklık kontrollü bir tünelde ve rijit montajda elektronik sürüklenme
muhtemelen baskındır: 0.1~$\mu$m/saat şartı bile günde $\sim2~\mu$m eder. Bu, Adım~6'da varsayılan
0.2~$\mu$m/$\sqrt{\text{gün}}$'den çok büyüktür. Önlem: pilot ton kalibrasyonu (bilinen bir referans
sinyaliyle elektroniğin sürekli kalibre edilmesi; SLS~2.0 ve diğer yeni ışık kaynaklarında kullanılır) ve
daha sık BBA. Soru 32'ye göre iki demetli rijit mimari $\sigma_\delta\approx1~\mu$m'yi tolere ediyor; bu,
0.1~$\mu$m/saat elektronikle $\sim10$ saatte bir BBA demektir.

%======================================================================
\section{T2: Rijit BPM (ana rapor 8.2)}
%======================================================================

\soru{28}{Yeniden kurma hatası $\sigma_\delta$ ile büyüdüğü halde sahte EDM neden büyümüyor?}

Çünkü büyüyen hata, spinin göremediği doğrultulardadır. $\sigma_\delta=1~\mu$m'de iki demetin kalan simetrik
artığının gücünün \%93'ü $m=0$ ve $m=1$'de~\BEL. Bu desenler halkanın tamamının ötelenmesi ve eğilmesidir;
demet bunları izler, kuadrupol merkezlerinden geçer, alan görmez. Spinin en duyarlı olduğu $m=2$--3'teki
artık $\sigma_\delta$ ile büyür ama çok küçük kalır ($\sigma_\delta=1~\mu$m'de rms $\sim0.16~\mu$m; soru 32)
ve sahte EDM sayısal tabanın altında kalır~\YOR.

\sekil{ek_sekil_s12_t2.pdf}{T2 (rijit BPM). Sol: sahte EDM, BPM--kuadrupol göreli kaymasına karşı. Sağ: $\sigma_\delta=1~\mu$m'de kalan simetrik artığın harmonik dağılımı; iki demetin artığı spinin duyarsız olduğu $m=0$--1'de.}
\soru{29}{Kalan artık neden $m=0,1$'de toplanıyor? Şans mı, kestiricinin yapısı mı? Spin $m=0$ ve $m=1$'e
neden duyarsız?}

\textbf{Yapısaldır, şans değil~\YOR.} Rijit BPM'de okuma $(R-I)d$'dir. Halkanın tamamı düzgün olarak $d_0$
kadar ötelenirse kapalı yörünge de aynen ötelenir: $R\,d_0=d_0$, yani $(R-I)d_0=0$. Bu desen okumada hiç
görünmez; $(R-I)$'nin sıfır doğrultusudur. Test paketi bunu doğrular: düşey düzlemde düzgün ötelemenin
okuması $0.056$, yani $(R-I)$ neredeyse tekil~\BEL. Eğim ($m=1$) için benzer bir yakın-değişmezlik vardır.
Kestirici bu doğrultuları göremez; $\delta$'nın bu doğrultulara düşen kısmı, $1/s_{\rm min}$ ile büyütülür.
Bu yüzden artık orada toplanır. Görülmeyen ve büyütülen aynı doğrultulardır.

\textbf{Spin neden duyarsız~\YOR.} Aynı nedenle: halka bir bütün olarak öteleniyorsa demet de onunla öteleniyor
ve her kuadrupolün merkezinden geçiyor. Kuadrupol dipol alanı görmeyen demette spin dönmez. Ölçülen
duyarlılık $m=0$'da $m=2$'nin $\sim10^{-4}$'ü (ana rapor Adım~5)~\BEL. Fiziksel olarak simetrik bir
düzenlilik: ``rijit BPM'in göremediği şey'' ile ``spinin umursamadığı şey'' aynıdır, çünkü ikisi de
demetin kuadrupole göre konumuna bakar. Bu, rijit BPM mimarisinin temel güzelliğidir.

\soru{30}{Üstünlüğün $\sigma_\delta$ arttıkça büyümesi ilk bakışta ters geliyor. Nedeni? ``Sayısal tabanla
sınırlı'' ne demek; gerçek üstünlük daha mı büyük?}

Üstünlük $=$ (tek demetin $C$'si) / (iki demetin $C$'si). Tek demette $\delta$ spinin duyarlı olduğu
$m=2$--3'e de düşer (\%37) ve $C\propto\sigma_\delta^2$ büyür: 0.034, 0.16, 8.2~\BEL. İki demette ise $\delta$
neredeyse yalnız $m=0$--1'e düşer (soru 29) ve $C$ sayısal tabanda kalır: 0.008, 0.007, 0.010~\BEL. Pay
büyürken payda sabit kaldığı için oran büyür.

``Sayısal tabanla sınırlı'': ölçüm aracı kaçıklık sıfırken bile $\approx0.008\times$hedef verir (kısaltılmış
ayarın sayısal gürültüsü). İki demetin $C$'si bu tabanın altına inemez görünür; gerçek değer daha küçük
olabilir. Yani 800 kat bir \emph{alt sınırdır}; gerçek üstünlük daha büyük olabilir~\YOR. Daha kesin bir
değer için tam ayar (tabanı daha düşük) ve daha uzun izleme gerekir. Ama bu pratik açıdan önemli değildir:
iki demet zaten hedefin 100 kat altındadır.

\soru{31}{$\sigma_\delta\lesssim0.3$--1~$\mu$m gereksinimi gerçekçi mi? Literatürde BPM--kuadrupol göreli
kararlılığı için ulaşılmış değerler neler?}

\LIT~Bulduğum değerler (hepsi elektron halkaları ya da lineer hızlandırıcılardan; doğrudan
``BPM--kuadrupol göreli, haftalar boyunca'' ölçen bir kaynak bulamadım):
\begin{itemize}[nosep]
\item SPring-8 yükseltmesi: komşu kuadrupole göre sıkı tolerans, BPM için 5~$\mu$m uzun süreli kararlılık
şartı; yeni elektronik prototipi uzun sürede 1~$\mu$m içinde~\cite{spring8}.
\item APS MBA: 7 günde $\le1~\mu$m sürüklenme şartı (ekleme bölgelerinde)~\cite{aps_mba}.
\item SLS~2.0: elektronik sürüklenme şartı 0.1~$\mu$m/saat~\cite{sls2}.
\item SRRC: BPM'in girder'a göre konumu bir vardiyada 1~$\mu$m içinde, sonra 0.1~$\mu$m'ye
iyileştirilmiş~\cite{lep}.
\item CLIC/PACMAN: BPM ve kuadrupol ortak destekte, elektromanyetik merkezler arası mesafe mikrometre
aralığında bulunmuş (ön hizalama, zamanla kararlılık değil)~\cite{pacman}.
\item TTF2 (DESY): kuadrupol içindeki BPM'ler gerili telle kuadrupol ekseninde kalibre edilmiş~\cite{ttf2}.
\end{itemize}
\sekil{ek_sekil_s15_bpm_lit.pdf}{BPM sürüklenmesi: literatür değerleri ve bu çalışmanın sayıları (süreler farklı; kaba karşılaştırma).}
\YOR~Sonuç: ``günler boyunca 1~$\mu$m'' bugünün iyi ışık kaynaklarının şartlarıyla aynı düzeydedir;
``0.3~$\mu$m'' zorlayıcıdır. Yeni sonuçlara göre (iki demet $\sigma_\delta=1~\mu$m'de hedefin 100 kat altında)
gereksinim $\sigma_\delta\lesssim1~\mu$m olarak yazılabilir; bu, ulaşılmış değerlerle uyumludur.

\soru{32}{Daha büyük $\sigma_\delta$'da ne olmasını bekliyoruz? Test edilmesi gerekir mi?}

\HIP~$C$, kalan artığın spin-duyarlı kısmının karesiyle ölçeklenir. İki demette o kısım $\sigma_\delta$ ile
doğrusal büyüdüğünden $C\propto\sigma_\delta^2$ beklenir. $\sigma_\delta=1~\mu$m'de $C\le0.01$ (taban; gerçek
değer daha küçük olabilir) olduğundan $C<1$ için $\sigma_\delta\lesssim10~\mu$m yeterli olmalıdır. Bu
kaba bir üst sınırdır, çünkü 1~$\mu$m'deki gerçek $C$ bilinmiyor; tabanın altındaysa sınır daha da gevşer.
Ama bu tahmin, ``$\delta$ BPM'den BPM'e bağımsız'' modeline dayanır. Gerçekte $\delta$ korelasyonlu olabilir
(ör.\ bütün elektroniğin ortak sıcaklıkla kayması); o zaman $\delta$'nın deseni farklı harmoniklere düşer.
Test edilmesi önerilir: $\sigma_\delta=3$ ve 10~$\mu$m, ayrıca korelasyonlu $\delta$ (birkaç saatlik iş)~\YOR.
Makalede bir gereksinim olarak yazılacaksa test edilmiş aralıkla sınırlı tutulmalı ($\le1~\mu$m).

%======================================================================
\section{Adım 5: Harmonikler (ana rapor 9)}
%======================================================================

\soru{33}{``$k=24$ bileşeni yaklaşık 100 kat bastırılır'' ifadesi ölçülen kazançlarla neden uyuşmuyor?
Analitik $G_k$ formülü nerede yanlış? Makalede nasıl düzeltmeliyiz?}

\textbf{Makalenin akıl yürütmesi.} Simetrik desende QF ve QD aynı yöne kayar. Gradyanları zıt işaretli
olduğu için tekmeleri zıt işaretlidir: $\theta_{\rm QF}=+g\,d$, $\theta_{\rm QD}=-g\,d$. Tekme deseni hücre
içinde $+,-,+,-$ biçiminde alternatif olur, yani 24. harmoniktir. Pürüzsüz odaklama yaklaşımında yörünge
tepkisi $G_k\approx C/|Q^2-k^2|$; $k=24$ için $\approx24.8/571\approx0.04$, düşük $k$ için $\sim$6. Oran
$\sim$100.

\textbf{Eksik varsayım~\YOR.} Bir tekmenin kapalı yörüngeye etkisi $\sqrt{\beta}\,\theta$ ile ağırlıklanır
($\beta$, tekme noktasındaki betatron fonksiyonu). FODO'da $\beta$, QF ve QD'de çok farklıdır:
\begin{center}\small
\begin{tabular}{lccc}
\toprule
düzlem & $\beta_{\rm QF}$ [m] & $\beta_{\rm QD}$ [m] & $(\sqrt{\beta_F}-\sqrt{\beta_D})/(\sqrt{\beta_F}+\sqrt{\beta_D})$ \\
\midrule
x & 65.3 & 35.5 & 0.15 \\
y & 41.5 & 75.8 & $-0.15$ \\
\bottomrule
\end{tabular}
\end{center}
(\kod{drift\_monitor/fodo\_lattice.py}, kuadrupol girişinde~\BEL.) Ağırlıklı tekmeler birbirini tam
götürmez: QF ve QD tekmeleri eşit ve zıt olsa da $\sqrt{\beta}$ farkı yüzünden hücre başına net bir tekme
kalır. Bu net tekmenin, antisimetrik desenin (tekmeleri aynı işaretli, $\sqrt\beta$'lar toplanır) net
tekmesine oranı yukarıdaki son sütundur: $\approx0.15$. Bu net tekme, düşük harmoniklerde
\emph{rezonant} biçimde yörünge üretir. Ölçülen simetrik/antisimetrik kazanç oranı düşük $m$'de 0.15--0.25
(ana rapor Tablo 11)~\BEL; basit $\sqrt\beta$ hesabıyla uyumlu.

\sekil{ek_sekil_s13_beta.pdf}{Sol: simetrik desende QF ve QD tekmelerinin $\sqrt\beta$ ile ağırlıklanmış değerleri; net tekme $\approx\pm0.15$. Sağ: ölçülen simetrik/antisimetrik yörünge kazancı oranı; düşük $m$'de $\sqrt\beta$ tahminiyle uyumlu, makalenin 100 kat iddiasının çok üstünde.}
Yani hata formülün kendisinde değil, simetrik desenin ``saf $k=24$'' sayılmasındadır: $\beta$ modülasyonu
simetrik desenin \%15'ini düşük harmoniklere sızdırır ve o kısım bastırılmaz.

\textbf{Makalede~\YOR.} (1)~Sec.~III'teki 100 kat bastırma iddiası kaldırılmalı. (2)~Yerine ``simetrik
desen, QF ve QD'deki $\beta$ farkı yüzünden $(\sqrt{\beta_F}-\sqrt{\beta_D})/(\sqrt{\beta_F}+\sqrt{\beta_D})
\approx0.15$ oranında düşük harmoniklere sızar; bu yüzden yörünge görünürlüğü antisimetrikten yalnızca
4--7 kat düşüktür'' gibi nicel bir ifade konmalı. (3)~Şekil 1(a)'daki analitik $G_k$ eğrisi ölçülen
harmonik kazançlarla değiştirilmeli ya da yanına konmalı. Bu düzeltme makalenin ana fikrini bozmaz:
simetrik bileşen yine \emph{daha az görünür}dür, yalnız ``kör'' değildir.

\soru{34}{``Simetrik bileşen yörüngedeki görünürlüğüne göre spine daha duyarlıdır, ama fark ılımlıdır.''
Bunun ana iddia için anlamı ne? Lehimize mi, aleyhimize mi?}

Karışık~\YOR.
\begin{itemize}[nosep]
\item \emph{Aleyhimize:} Yörüngede 1 birim görünen simetrik desen, 1 birim görünen antisimetrik desenden
2--3 kat fazla sahte EDM üretir. Yani simetrik bileşeni yörüngeden kestirmek, aynı yörünge hassasiyetiyle
antisimetrikten daha zor bir hedefi karşılamak demektir.
\item \emph{Lehimize:} Fark 2--3 kattır, ``50 kat'' değil (ilk tahmin yanlıştı). Simetrik bileşen yörüngeye
kör değil (soru 33), spine aşırı duyarlı da değil. Yani yörünge tabanlı bir yöntemin onu yakalaması
fiziksel olarak mümkündür; mesele yalnızca 4--7 kat daha iyi yörünge hassasiyetidir. İki demet bunu
sağlıyor.
\item \emph{Not:} Adım~5'in sahte EDM sayıları eski prosedürle ölçüldü; oranlar aynı desenler arasında
olduğundan muhtemelen daha az etkilenir, ama doğrulanmadı.
\end{itemize}

%======================================================================
\section{Adım 6: Zaman senaryosu (ana rapor 10)}
%======================================================================

\soru{35}{Gürültü modeli gerçekçi mi? 10~$\mu$m/tur ve 10~nm taban nereden geliyor? Tur sayısıyla ortalama
alma gerçekten $\sqrt N$ kazanç verir mi?}

\textbf{Değerler~\LIT.} FCC-ee çalışması tur başına 10~$\mu$m ve yörünge için 1~$\mu$m'den iyi çözünürlük
hedefler~\cite{fccee}. NSLS-II tek demet paketi için tur başına mikrometre çözünürlüğe ulaşmıştır; uzun paket
trenlerinde mikrometre altı rutindir~\cite{nsls2}. HEPS düğme BPM'leri 10~Hz'de 0.1~$\mu$m'den,
220~kHz'de 1~$\mu$m'den iyi çözünürlük belirtir~\cite{heps}. Yani 10~$\mu$m/tur muhafazakâr bir değerdir;
ama bunlar elektron halkalarıdır. Proton demeti daha uzun paketlidir ve sinyal biçimi farklıdır; proton
EDM halkası için doğrudan bir değer bulamadım.

\textbf{$\sqrt N$ kazanç~\YOR.} Yalnız birbirinden bağımsız (beyaz) gürültü için geçerlidir. Gerçek BPM
gürültüsünde düşük frekanslı bileşenler (elektroniğin ısıl sürüklenmesi, $1/f$ gürültüsü) vardır ve
ortalama bir noktadan sonra iyileşmez. Modeldeki ``taban'' terimi tam olarak bunu temsil eder: ortalama
uzadıkça $\sigma_n$, beyaz kısım sıfıra gitse bile tabana iner. 60~s ortalamada beyaz kısım yalnızca
$\sim2.7$~nm'dir (224~kHz ile); yani sonuç neredeyse tamamen 10~nm tabanla belirlenir. Bu taban da
\HIP'tir: HEPS'in 10~Hz'de $<0.1~\mu$m değeri, 60~s'de 10~nm'nin ulaşılabilir olduğunu düşündürür, ama
günler boyunca bu düzey elektroniğin sürüklenmesiyle sınırlıdır (soru 27). Önemli ayrım: kısa süreli
taban (ölçüm gürültüsü) ile uzun süreli sürüklenme ($\delta$) farklı şeylerdir; ikincisi T2 ve Adım~6'da
ayrıca modellendi.

\textbf{Küçük düzeltme~\BEL.} Adım~6'da tur frekansı 270~kHz alınmıştı; doğrusu 224~kHz. $\sigma_n$
10.31 yerine 10.37~nm olur; sonuçlara etkisi ihmal edilebilir. Tespit sınırı örneğinde ($T_{\rm ort}=1$~s)
beyaz kısım 19.2 yerine 21.1~nm, sınır 49 yerine 52~nm olur.

\soru{36}{``İki demet bir hafta boyunca 25 ölçümün hiçbirinde 0.18'i aşmıyor.'' 5 seed ile bu sonuca ne
kadar güvenebiliriz?}

\YOR~İki ayrı soru var. (1)~\emph{Hedefi aşma olasılığı:} 25 ölçümde hiç aşma yoksa, aşma olasılığı için
\%95 üst sınır kabaca $3/25=\%12$'dir (``üç kuralı''). Ama 25 ölçüm bağımsız değildir (aynı seed'in farklı
saatleri), etkin örneklem daha küçüktür. Bu yüzden ``olasılık \%12'den küçük'' bile iyimserdir. (2)~\emph{Marj:}
Daha güçlü argüman büyüklüktür. En büyük değer 0.18; hedef 1. Seed'ler arası dağılım yaklaşık lognormal ve
genişliği $\sim$3 kat (16--84\% aralığı) ise 1'e ulaşmak $\sim$2 standart sapma uzakta bir kuyruk ister.
\HIP: hedefi aşma olasılığı bir yüzde mertebesindedir. Makaleye girecek bir sayı için 20--50 seed
önerilir; bu, kısaltılmış ayarla yaklaşık bir gün sürer.

\sekil{donmus_adim6.pdf}{Ana rapor Şekil 7 (tekrar): bir haftalık senaryo. Orta panel: sahte EDM, 5 seed; iki demetin (koyu turuncu) bütün noktaları 0.18'in altında.}
\soru{37}{Gevşetilmiş BBA artığı gereksinimi ne olmalı ve hangi sonuca dayanıyor?}

\YOR~Dayanak: $t=0$ (yalnız 0.3~$\mu$m BBA artığı) için $C$ medyan 0.063, en büyük 0.14 (5 seed)~\BEL. Sahte
EDM $\sigma^2$ ile ölçeklendiği için ($p=2.0002$)~\BEL, BBA artığı $\sigma_{\rm BBA}$ için
$C\approx0.14\,(\sigma_{\rm BBA}/0.3~\mu\text{m})^2$ (en kötü seed). $C<1$ için $\sigma_{\rm BBA}<0.8~\mu$m.
İki demetli izleme BBA artığının üstüne bir miktar ekler (hafta boyunca en büyük 0.18), bu yüzden 2--3 kat
güvenlik payıyla: \textbf{$\sigma_{\rm BBA}\lesssim0.4$--0.5~$\mu$m}. Eski gereksinim (0.15--0.2~$\mu$m)
eski prosedürün yanlı değerlerinden geliyordu. Bu ölçekleme 5 seed'e dayanır~\HIP.

\sekil[0.7]{ek_sekil_s14_bba.pdf}{BBA artığı gereksinimi: $t=0$'daki ölçülen değerler ve $\sigma^2$ ölçeklemesi.}
\soru{38}{``Makale modeli hedefte marjsız.'' Makalede bunu nasıl ifade etmeliyiz? İyi haber mi, kötü haber mi?}

\YOR~İkisi birden. \emph{İyi:} BPM'ler zeminle hareket etse bile makalenin yöntemi felaket değil; sahte
EDM'yi ham değerden (55) hedefin çevresine (0.6--1.3) indiriyor. Yani yöntem bu koşulda da \emph{büyük bir
iyileştirme} sağlıyor. \emph{Kötü:} Hedefin çevresinde olmak, bir EDM deneyi için yeterli değildir; 5 seed'in
4'ünde 96. saatte hedef aşılıyor. Makale bu koşulda tek başına yeterli olduğunu iddia edemez.

Makalede önerilen anlatım: BPM modeli bir \emph{senaryo parametresi} olarak sunulmalı. Lab-sabit BPM'ler
sabit kalıyorsa makale modeli yeterlidir; BPM'ler yerle hareket ediyorsa model, BPM hareketini içerecek
biçimde genişletilmelidir (rijit montaj + $\rho$-bilen kestirici), aksi halde marj kalmaz. Bu, okuyucuya
yöntemin geçerlilik alanını dürüstçe gösterir ve hakemin kaygısına doğrudan cevaptır.

%======================================================================
\section{Makale için çıkarımlar (ana rapor 11)}
%======================================================================

\soru{39}{Daralmış iddia makalenin özgünlüğünü ve değerini ne kadar etkiliyor? Hakemlere nasıl sunmalıyız?}

\YOR~Özgün kısım, ``simetrik bileşen iki demetin toplam kanalında görünür ve izlenebilir'' fikridir. Bu
fikir kontrollerden sağ çıktı: doğru koşullarda iki demet simetrik sürüklenmeyi tek demetten çok daha iyi
izliyor. Daralan şey fikrin kendisi değil, \emph{koşulsuz} sunuluşudur.

Hatta daralmış biçim, bazı yönlerden daha güçlü bir katkıdır: (1)~rijit BPM ile iki demetin
$\delta$'ya karşı sağlamlığı (soru 29), yapısal ve fiziksel bir açıklaması olan yeni bir sonuçtur;
(2)~geçerlilik alanını ve tasarım gereksinimlerini nicel olarak veren bir makale, bir deney tasarımı
için doğrudan kullanılabilir.

Hakemlere sunarken: Hakem A'nın üç kaygısının (saptırıcı hataları, korelasyonlu hatalar, uzun dalgalı
hareket) her biri ayrı bir simülasyonla test edildi ve sonuçlar --- olumsuz olanlar dahil --- makaleye
girdi. Bu, ``iddiayı daralttık'' değil, ``iddianın geçerli olduğu koşulları gösterdik'' diye çerçevelenebilir.
Olumsuz sonuçları açıkça yazmak (lab-sabit BPM'de makale modelinin marjsız kalması, saptırıcı kararlılığının
zorluğu) hakemlerin güvenini artırır.

\soru{40}{Tasarım gereksinimleri tablosunun tamamını tek tek açıklar mısın? Her gereksinim ne kadar
gerçekçi, hangisi en zor?}

\begin{center}\small
\begin{tabularx}{\textwidth}{>{\raggedright\arraybackslash}p{4.2cm}>{\raggedright\arraybackslash}X>{\raggedright\arraybackslash}p{3.3cm}}
\toprule
gereksinim & ne demek, neden & gerçekçilik \\
\midrule
BPM--kuadrupol göreli kararlılığı $\sigma_\delta\lesssim1~\mu$m (BBA aralığı boyunca) & BBA'dan sonra BPM
merkezi ile kuadrupol merkezi arasındaki mesafe bu kadar değişmemeli (T2) & ulaşılmış düzeyde; elektronik
sürüklenme belirleyici (soru 27, 31) \\
\addlinespace
kuadrupol üstü manyetik dipol trim sargıları & yalnız bu düzeltici türü kaçıklığı yerinde iptal eder
(Adım 3, soru 21) & standart teknoloji \\
\addlinespace
saptırıcı başına yüksek gerilim ve düşey elektrik trimi; saptırıcı sütunlu fit & saptırıcı sürüklenmesini
kestirip kaynağında geri almak (T1) & trimin kendisi kolay (soru 25) \\
\addlinespace
\textbf{saptırıcı başına alan kararlılığı $\sim3\times10^{-7}$ (bağımsız, yavaş)} & T1 bu düzeyde test edildi;
$10^{-5}$'te $C\sim10$'a çıkması beklenir (soru 25) & \textbf{en zor}: bugünkü kaynaklar
$10^{-5}$--$10^{-4}$ \\
\addlinespace
yalnız kuadrupol fiti kullanılırsa saptırıcı kararlılığı $\lesssim5\times10^{-8}$ ve 0.05~$\mu$rad & saptırıcı
sütunu olmayan fitin sızıntısı $\sim1~\mu$m / ($10^{-7}$) (Adım 2) & gerçekçi değil; saptırıcı sütunları
şart \\
\addlinespace
BBA artığı $\lesssim0.4$--0.5~$\mu$m & izleyici BBA kalitesini korur, daha iyisini yapamaz (soru 37) &
ulaşılmış düzeyde~\HIP \\
\addlinespace
BPM elektronik tabanı $\lesssim10$--100~nm & tespit sınırını belirler (Adım 6) & kısa sürede ulaşılmış;
günler boyunca zor (soru 35) \\
\addlinespace
çift kısmın sadeleşmesi için CW/CCW simetrisi & roll gibi birinci mertebe etkiler yalnız farkta sadeleşir
(soru 20) & zor; bağıl simetri $\sim10^{-8}$ gerekebilir (tahmin) \\
\bottomrule
\end{tabularx}
\end{center}
En zor iki gereksinim saptırıcı alan kararlılığı ve CW/CCW simetrisidir. İkisi de bu yönteme özgü değildir;
her proton EDM halkasının genel sorunlarıdır. Ama yöntem, saptırıcı sürüklenmesine özellikle duyarlıdır,
çünkü saptırıcı hataları simetrik sinyalle aynı kanala düşer~\YOR.

\soru{41}{Sec.~V.B'deki ``ortak hareket'' argümanının tersine çevrilmesini daha ayrıntılı açıklar mısın?
Makalede bu kısmı nasıl yazmalıyız?}

\textbf{Makalenin bugünkü argümanı~\YOR.} ``BPM'ler kuadrupollerle birlikte hareket eder, bu yüzden ortak
hareket yöntemi etkilemez'' (BPM'in bir kusuru olarak değil, zararsız bir ayrıntı olarak). Ama
simülasyonlarda BPM'ler lab-sabit modellenmişti (Adım~0). Ortak hareket gerçekte modelde olsaydı, makale
modeli (E3) BPM'lerin hareketini kuadrupol kaymasıyla karıştırırdı: $\rho=1$'de simetrik sürüklenmenin
yalnız küçük kısmını bulur (kalan 4.4~$\mu$m; sahte EDM 0.19, marjsız)~\BEL.

\textbf{Tersine çevirme.} Ortak hareket bir sorun değil, bir \emph{kaynaktır}, ama yalnız kestirici onu
biliyorsa. BPM kuadrupole rijit bağlıysa okuma, demetin kuadrupol merkezine göre konumudur; bu, sahte EDM'yi
belirleyen büyüklüğün ta kendisidir (soru 13). Bunu bilen iki demetli kestirici (E4) hem simetrik
sürüklenmeyi bulur hem de BPM--kuadrupol göreli kaymasına karşı yapısal olarak sağlamdır (soru 29).

\textbf{Makalede nasıl~\YOR.} (1)~BPM modelini açıkça bir denklemle tanımla: $\Delta r_b=R_b\Delta d-u+n_b$,
$u=\rho\Delta d$. (2)~Lab-sabit ($\rho=0$) ve rijit ($\rho=1$) durumları ayrı alt başlıklarda ver.
(3)~Rijit durumda kestiricinin $(R-\rho I)$ kullanması gerektiğini ve bunun neden doğal olduğunu
(okuma $=$ demet $-$ kuadrupol merkezi) anlat. (4)~Rijit montajın lineer hızlandırıcılardaki örneklerini
(TTF2/XFEL, CLIC/PACMAN) ve proton halkasında senkrotron ışıması olmadığı için ışık kaynaklarındaki
karşı-argümanın geçmediğini yaz. (5)~$\rho$'nun nasıl bilineceğini (tasarım ve/veya ölçüm; soru 15) ve
yanlış bilinirse ne olacağını ver.

\soru{42}{Düzenlileştirme makalede nasıl anlatılmalı? \kod{rcond} seçiminin gerekçesi nasıl verilmeli?}

\YOR~(1)~Yöntemin bir \emph{kesik tekil değer ayrışımı} (TSVD) olduğunu ve eşiğin anlamını (tekil değer
oranı) yaz. (2)~Eşiğin keyfi seçilmediğini göster: gerçek kaçıklığı bilmeyen standart bir kural
(uyuşmazlık ilkesi: fit artığı gürültü düzeyine inene kadar en çok düzenlileştir) kullan ve bunun ``kahin''
seçime çok yakın sonuç verdiğini bir cümleyle belirt (soru 14). (3)~Tek ve iki demeti \emph{aynı kuralla}
karşılaştır; sabit 0.01 eşiği tek demeti haksız biçimde kötü gösterir (soru 12). (4)~Bonus: bu kural
modelin yanlış olduğunu (fitin gürültü düzeyine inememesi) kendiliğinden gösterir; bu bir operasyonel
tanı aracı olarak sunulabilir.

\soru{43}{Sıfır ofset varsayımını makalede nasıl savunuruz?}

\YOR~Savunmak yerine yeniden tanımlamak daha doğru (soru 5, 11): ``Yöntem BBA referansına göre değişimleri
ölçer. Statik BPM ofsetleri referansta sadeleşir. Geriye kalan, ofsetin ve BPM--kuadrupol mesafesinin
referanstan sonraki değişimi $\delta$'dır; bunun etkisi Sec.~X'te nicel olarak verilmiştir.'' Böylece
Şekil~4'ün koşulu (``sıfır ofset'') aslında ``referanstan beri $\delta=0$'' olur ve $\delta\ne0$ ayrıca
gösterilir. Mutlak ofsetle (referans olmadan) yöntemin çalışmadığı da bir cümleyle söylenmelidir; bu,
okuyucunun yöntemi yanlış kullanmasını önler.

\soru{44}{Hangi sayılar etkileniyor ve yeniden ölçüm ne kadar iş? (Madde 9'daki itirazımla birlikte
düşünelim.)}

\textbf{Soru 9'un sonucuyla birlikte~\YOR.} Fırlatma noktasına bağımlılığın nedeni değişti (potansiyel
enerji değil, sert kenarlı saptırıcılarda hücreden hücreye $|p|$ farkı), ama sonuç değişmedi: eski
prosedürle ölçülen sahte EDM, seçilen enerjiye bağlı ve yanlıdır; donmuş spin prosedürü bu yanlılıktan
arınmıştır. Bu yüzden makaledeki bütün sahte EDM sayıları etkilenir.

\textbf{Etkilenen sayılar~\HIP} (makale metnini burada yeniden incelemedim; ana raporda kullanılan kurulumlara
göre): Şekil~4'ün üç değeri (ham, tek demet, iki demet) ve seed noktaları; $\sigma^2$ ölçekleme şekli ya da
metindeki ilgili değerler; varsa harmonik/tune taramalarındaki sahte EDM değerleri. Doğrusal (yörünge,
yeniden kurma) sonuçlar etkilenmez.

\textbf{İş~\BEL.} Kısaltılmış ayarla bir sahte EDM ölçümü (bir yön) 4 çekirdekte $\sim$1.5 dakikalık iş
hacmine karşılık gelir (4 paralel ~6 dakikada 4 ölçüm). Şekil~4 (3 durum $\times$ 15 seed $\times$ 2 yön
$=90$ ölçüm) $\approx2.5$ saat; $\sigma^2$ (3 $\sigma$ $\times$ 5 seed $\times$ 2 yön) $\approx1$ saat. Şekil~4
ve $\sigma^2$ için zaten donmuş spin sonuçları var (5 ve 3 seed); 15 seed'e tamamlamak birkaç saatlik iştir.

%======================================================================
\section{Sınırlamalar (ana rapor 12)}
%======================================================================

\soru{45}{Statik başlangıç durumu ne kadar büyük bir açık? Simüle etmek için ne gerekir? Hakemler sorarsa
ne deriz?}

\textbf{Açığın büyüklüğü~\YOR.} Büyük, ama bu yönteme özgü değil. $10^{-4}$/0.1~mrad statik saptırıcı
hataları milimetre düzeyinde kapalı yörünge yaratır (Adım~2: tek saptırıcıdan 418~$\mu$m) ve roll'un birinci
mertebe etkisi $3\times10^{-7}$'deki $1.26\times10^{8}$'den 300 kat büyük olur (soru 20). Bu durumdan,
herhangi bir ölçüme başlamadan önce, kalibrasyonla çıkılmalıdır: statik yörünge düzeltmesi, BBA,
saptırıcıların hizalanması. Bizim senaryomuz bu kalibrasyonun bittiği anda başlar (``BBA sonrası''); kalibrasyonun
kendisinin bıraktığı artık, $\sigma_{\rm BBA}$ ve kalan statik saptırıcı hataları olarak senaryoya girer.

\textbf{Simüle etmek için~\HIP.} (1)~Statik saptırıcı ve kuadrupol hatalarıyla başla (mm yörünge). (2)~Klasik
yörünge düzeltmesi ve BBA'yı simüle et (gürültülü, yinelemeli). (3)~Kalan durumu bizim senaryomuzun
başlangıcı yap. Zorluklar: mm düzeyinde doğrusal olmayan yörünge etkileri, büyük sahte EDM'nin donmuş spin
aramasını zorlaması, ve kalibrasyon prosedürünün kendisinin (hangi BBA yöntemi) bir tasarım kararı olması.
Birkaç günlük iş.

\textbf{Hakemlere~\YOR.} ``Makale, BBA sonrası sürüklenmenin izlenmesini ele alır. Statik hataların
kalibrasyonu ayrı ve standart bir problemdir; bizim sonuçlarımız kalibrasyon kalitesi cinsinden
($\sigma_{\rm BBA}\lesssim0.5~\mu$m, kalan saptırıcı hataları) parametrize edilmiştir.'' Bu dürüst ve
savunulabilir bir çerçevedir; ama hakem bunun ulaşılabilirliğini sorarsa, Omarov vd.\ ve proton EDM
fizibilite çalışmalarının ilgili bölümlerine atıf gerekir~\cite{omarov,cpedm_feas}.

\soru{46}{Adım 6'da saptırıcı sürüklenmesi ve T1 trimi birleştirilirse sonucun değişmesini bekliyor muyuz?
Yapılması gerekir mi?}

\HIP~$3\times10^{-7}$ düzeyinde saptırıcı sürüklenmesi ve trim ile T1 tek başına $C\approx0.01$--0.06
veriyor; Adım~6'nın iki demetli sonucu 0.05--0.18. İki katkı büyük ölçüde bağımsız ise toplam
$\approx0.05$--0.2 olur; yani sonuç değişmez. Ama iki önemli soru var: (1)~saptırıcı sürüklenmesinin
\emph{zaman yasası} ne (yer hareketi gibi $\sqrt t$ mi, sıcaklıkla mı, doğrusal mı)? (2)~soru 25'e göre
gerçekçi kararlılık $10^{-5}$ ise sonuç $C\sim10$ olur. Yani birleştirmenin asıl değeri, ``hangi saptırıcı
kararlılığında yöntem hâlâ çalışır'' sorusuna nicel cevap vermesidir. Önerilir: Adım~6 + saptırıcı
sürüklenmesi, kararlılık taraması $10^{-7}$--$10^{-5}$ (bir günlük iş).

\soru{47}{Adım 5'in sahte EDM sayıları makalede kullanılacak mı? Kullanılacaksa yeniden ölçülmeleri şart mı?}

\YOR~Makalede yer alacaksa evet, şart. Eski prosedürün hatası desene bağlı (soru 10'daki oran seed'e göre
$-4.6$ ile $+2.4$ arasında) ve harmonik desenler çok farklı yörüngeler ürettiği için harmonikler arasındaki
oranlar da etkilenebilir. Adım~5 için 13 harmonik $\times$ 2 tür $\times$ 2 faz $\times$ 2 yön
$\approx100$ ölçüm, kısaltılmış ayarla $\sim$3 saat. Hangi kısmın makalede kullanılacağına göre:
\begin{itemize}[nosep]
\item yörünge kazançları (doğrusal): olduğu gibi kullanılabilir, etkilenmez;
\item ``spin $m=0$'da çok bastırılır, $m=2$--3'te en duyarlı'' nitel ifadesi: büyük farklara dayandığı için
büyük olasılıkla doğru kalır, ama bir kontrol ölçümü (birkaç harmonik) önerilir;
\item tablo halinde $C/\mu$m$^2$ değerleri: yeniden ölçülmeli.
\end{itemize}

%======================================================================
\section*{Kaynaklar}
\addcontentsline{toc}{section}{Kaynaklar}
%======================================================================
\begin{thebibliography}{99}\small
\bibitem{jedi_phaselock} N.~Hempelmann et al.\ (JEDI), ``Phase locking the spin precession in a storage
ring'', Phys.\ Rev.\ Lett.\ 119, 014801 (2017). (Bilinen kaynak; bu oturumda web'den yeniden doğrulanmadı.)
\bibitem{omarov} Z.~Omarov et al., ``Comprehensive symmetric-hybrid ring design for a proton EDM experiment
at below $10^{-29}\,e\cdot$cm'', Phys.\ Rev.\ D 105, 032001 (2022). (Repoda PDF olarak mevcut.)
\bibitem{cpedm_feas} CPEDM / JEDI, ``Storage ring to search for electric dipole moments of charged particles
--- feasibility study'', \url{https://arxiv.org/pdf/1912.07881}.
\bibitem{fug} P.~Schmidt-Wellenburg, T.~Schmidt (PSI, FuG), ``The FuG HCB 200~kV bipolar high-voltage power
supply for the storage chamber in a nEDM experiment'',
\url{https://instrumentengineers.org/wp-content/uploads/2020/06/The-FUG-HCB-200kV-Bipolar-high-voltage-power-supply-for-storage-chamber-in-a-nEDM-experiment-Dr.-Schmidt-Wellenburg-Philipp-PSI-Schmidt-Thomas-FuG.pdf}.
\bibitem{hirfl} ``High-voltage detuning power system of HIRFL-CSRm electron cooler'', Nucl.\ Instrum.\
Meth.\ A (2022), \url{https://www.sciencedirect.com/science/article/abs/pii/S0168900222009913}.
\bibitem{sls2} ``Development of the SLS 2.0 BPM system'', IBIC2023, \url{https://proceedings.jacow.org/ibic2023/papers/mo3c03.pdf}.
\bibitem{aps_drift} ``Sensitivity and offset calibration for the beam position monitors at the Advanced
Photon Source'', \url{https://www.researchgate.net/publication/3634890}.
\bibitem{spring8} ``Development status of a stable BPM system for the SPring-8 upgrade'', IBIC2016,
\url{https://proceedings.jacow.org/ibic2016/papers/tupg06.pdf}.
\bibitem{aps_mba} R.~Lill, N.~Sereno, B.~Yang, ``BPM stability studies for the APS MBA upgrade'',
\url{https://www.osti.gov/servlets/purl/1393935}.
\bibitem{lep} ``Calibration of centre-of-mass energies at LEP 2 for a precise measurement of the W boson
mass'', \url{https://arxiv.org/pdf/hep-ex/0410026}; ve ``position monitors bpm'' konu derlemesi (SRRC
girder kararlılığı), \url{https://www.science.gov/topicpages/p/position+monitors+bpm.html}.
\bibitem{pacman} CLIC/PACMAN ön hizalama tezi, Univ.\ Pisa, \url{https://etd.adm.unipi.it/t/etd-12172016-033103/}.
\bibitem{ttf2} TTF2 BPM--kuadrupol birimlerinin gerili tel test tezgâhı, LINAC2004,
\url{https://epaper.kek.jp/l04/PAPERS/TUP69.PDF}.
\bibitem{fccee} FCC Feasibility Study Report, Vol.~2, \url{https://arxiv.org/pdf/2505.00274}.
\bibitem{nsls2} ``Single micron single-bunch turn-by-turn BPM resolution achieved at NSLS-II'', IPAC2016,
\url{https://proceedings.jacow.org/ipac2016/papers/weobb01.pdf}.
\bibitem{heps} ``Beam position monitor characterization for the High Energy Photon Source synchrotron'',
Symmetry 15, 660 (2023), \url{https://doi.org/10.3390/sym15030660}.
\end{thebibliography}
Not: Web kaynaklarının değerleri arama özetlerinden alındı; makaleye girecek değerler için orijinal
belgelerin kontrol edilmesi önerilir.

\end{document}
